% Commentary by GPT-6 Astra (OpenAI, ChatGPT Pro mode).
% Chapter XVI, sections 1--3. Source checked against the attached PDF.

\begin{dcomment}{context}{p2}{Axioms for equivalences before a choice of model}
\textit{Commentary and mathematical audit by GPT-6 Astra (OpenAI, ChatGPT Pro mode).}

Grothendieck starts with a class of functors, rather than with cofibrations and fibrations. The distinction matters throughout this chapter: the latter classes are to be recovered, when possible, from stability properties of the former. His letter to Thomason of 2 April 1991, \S1, describes precisely this separation between the equivalences carrying the homotopy theory and the auxiliary devices used to calculate it \citep{Grothendieck1991Letter}. Here the basic datum is a class $W$ in $\Cat$; the localization $L:\Cat\to\Hot_W$ exists in the chosen larger universe and need not come with a calculus of single fractions.

Three meanings of saturation must consequently remain distinct. The two-out-of-three property controls a composable pair. The extra condition (1c) controls a pair $f:A\to B$, $g:B\to A$ when both composites are already equivalences. It is just what is needed to recognize a homotopy inverse before one knows that every arrow inverted by $L$ belongs to $W$. Closure under retracts is an additional requirement here. Strong saturation means $W=L^{-1}(\operatorname{Iso})$; it implies retract closure because an arrow which is a retract of an isomorphism is an isomorphism. One cannot substitute strong saturation silently in arguments made under (1).

There is also a consequential change of terminology in the later literature. Definition 1.1 calls the class satisfying (1,2,3) a fundamental localizer. Later authors generally call this a \emph{weak} fundamental localizer and reserve the unqualified name for the relative axiom (3 bis). The distinction concerns the hypothesis of Theorem A, not merely a word. In (3), the comma categories $X/s$ are individually aspherical. In (3 bis), the two comma categories may each have substantial homotopy, provided that the comparison between them is an equivalence. The latter implication does not follow by applying the former separately to $X\to S$ and $Y\to S$.

The editor's addition of Der (2) in footnote 1 is therefore mathematically substantive. A pointwise comparison of diagram objects gives an isomorphism only when the evaluations detect isomorphisms. The derivator axioms numbered here are those of Chapter I; importing a modern numbering without translating it would conceal the missing conservativity assumption. This is the same point already isolated in the commentary to XIIa on constructing localizers from derivators.
\end{dcomment}

\begin{dcomment}{exposition}{p4}{Comma categories replace a map by a fibration}
The construction $\Phi_S(X)$ makes Proposition 1.3 transparent. For $p:X\to S$, its objects can be written
\[
(s,x,\alpha:s\longrightarrow p(x)).
\]
A morphism to $(t,y,\beta)$ is a pair $(h:s\to t,k:x\to y)$ satisfying $p(k)\alpha=\beta h$. Projection to $s$ is a split Grothendieck fibration; reindexing along $h:s\to t$ replaces $(t,y,\beta)$ by $(s,y,\beta h)$. Thus its fibre over $s$ is $s\backslash X$. The inclusion
\[
i_X(x)=(p(x),x,\id_{p(x)})
\]
has left adjoint $q(s,x,\alpha)=x$: a map from $(s,x,\alpha)$ to $i_X(y)$ is determined uniquely by $k:x\to y$. In particular, $qi_X=\id_X$, and the unit supplies the other homotopy. This explains the comparison axiom (3$'$b), including the direction of its adjoint.

If $p:X\to S$ is itself a fibration, the inclusion $X_s\to s\backslash X$ has a right adjoint: choose the cartesian pullback $\alpha^*x$ of $x$ along $\alpha:s\to p(x)$. The adjunction follows from the defining universal property of a cartesian arrow. Its comma categories have terminal objects, so (2) and (3) put the inclusion in $W$. Consequently the fibrewise hypothesis can be transferred to the comma categories. This is the sole place where the fibration condition is doing work; a general functor's strict fibres need not compute the comma categories.

For example, take $S=(0\to1)$ and let $X=e$ map to $1$. The strict fibre over $0$ is empty, although $0\backslash X=e$. Replacing $X$ by $\Phi_S(X)$ produces $S$ itself. A condition on strict fibres of an arbitrary map would incorrectly reject this contractible example; the replacement repairs exactly that defect.

The reverse implication in the proposition first gets $\Delta^1\times X\to X$ from the fibre axiom and the segment axiom. The two endpoint inclusions become equal inverses after localization. A natural transformation therefore makes its endpoint functors equal in $\Hot_W$, and two-out-of-three gives the corresponding equivalence of membership in $W$. For a homotopy equivalence, (1c) then recognizes the two inverse functors as members of $W$. A category with a final or initial object is contractible by these natural transformations, which recovers (2). Finally the displayed square involving $\Phi_S(f)$ transfers the relative fibre condition back to $f$. This proof also explains why a chosen splitting and preservation of its cartesian arrows are unnecessary in the final fibrewise criterion: the comparison with the comma categories works for every functor over $S$ between fibrations. For the later systematic separation of fibrations and asphericity structures, see \citet[\S1]{Maltsiniotis2005Asphericite}.

Two slips in the printed proof should not change the argument. The last conclusion on ms.~5 is $i_X\in W$, rather than its printed repetition $i_X^\circ\in W^\circ$. On ms.~7 the fibres $s\backslash X$ are \emph{aspherical categories}; coasphericity properly describes the original map $X\to S$.
\end{dcomment}

\begin{dcomment}{later}{p8}{The weak minimum and the stronger axiom}
There is a definitive later answer to one foundational issue that must be distinguished from the stronger-axiom problem. Let
\[
W_\infty=\{f\mid |Nf|\text{ is a weak homotopy equivalence}\}.
\]
Theorem 6.1.18 of \citet[pp.~242--243]{Cisinski2006} proves that $W_\infty$ is the \emph{smallest weak fundamental localizer}. The cited pagination is that of the published Ast\'erisque volume; an earlier author version places the same theorem at pp.~213--214. Definition 6.1.12 and Proposition 6.1.16, pp.~240--241, are indispensable for checking that this theorem really concerns the weak axiom used here. His split-retraction saturation condition follows from (1c), while $W_\infty$ itself satisfies (1c), so the minimum comparison applies to the precise class of localizers in this manuscript. Thus every $W$ of Definition 1.1 contains $W_\infty$, and the intersection of all such $W$ equals $W_\infty$. It is not necessary to assume (3 bis) for this conclusion.

This does not prove that every weak localizer satisfies (3 bis). A theorem identifying the least objects in two partially ordered classes does not identify the two classes. Chiche explicitly records the latter problem as open in his 2014 thesis, Remark 2.1.15, and Cisinski still reported no example distinguishing the classes in his dated 2019 discussion \citep{Chiche2014,Cisinski2019WeakLocalizers}. We found no resolution in the targeted literature search for this commentary. These are dated statements of the problem's documented status, not a claim that a search can certify the absence of every subsequent result.

There is a further useful distinction. Strong fundamental localizers in the later sense are strongly saturated, by \citet[Proposition 4.2.4]{Cisinski2006}. That theorem may be used after (3 bis) has been obtained. It cannot supply strong saturation for an arbitrary $W$ satisfying just Definition 1.1. In particular, neither the finite-product arguments below nor the speculative reverse-excision axioms should acquire extra hypotheses invisibly through a change in terminology.
\end{dcomment}

\begin{dcomment}{exposition}{p9}{Sums and products require different hypotheses}
The two halves of the sum axiom do different jobs. Closure under sums builds a global equivalence from component equivalences. Pullback along a summand recovers the component equivalences from the global one. Indeed, in a coproduct map $\coprod X_i\to\coprod Y_i$, the inverse image of the summand $Y_i$ is exactly $X_i$. This proves Proposition 2.1. For finite sums, the factorization
\[
f\amalg g=(\id\amalg g)(f\amalg\id)
\]
proves that closure is equivalent to the condition that each $\emptyset\to X$ is an intrinsic $W$-cofibration. Empty and one-term sums cause no problem because identities belong to $W$.

The parenthesis in Corollary 2.2 should say Loc (4\emph{a}), rather than Loc (4). A coproduct is the filtered colimit of its finite partial coproducts, so filtered-colimit stability promotes finite \emph{closure} to arbitrary closure. This argument gives no summand-reflection statement (4b). The later summary on ms.~15$'$ correctly supplies a separate fibration axiom for that half. For an explicit counterexample, let $W_{\mathrm{ne}}$ consist of all functors between nonempty categories, together with $\id_\emptyset$. This is even a strong localizer: a relative comma comparison cannot belong to it when its source is empty and its target has a chosen object. It is closed under arbitrary coproducts and filtered colimits, since both constructions preserve the test of nonemptiness. However, $e\to e\amalg e$ is in $W_{\mathrm{ne}}$, while its pullback over the other summand is $\emptyset\to e$, which is not. Thus (4b) really fails. The printed parenthesis is retained as an edition slip.

For products the editors' qualification is equally necessary. With (3 bis), consider $f\times\id_C:A\times C\to B\times C$ over $C$. Over $c$ the comparison is $f\times\id_{C/c}$. Since $C/c$ has a final object, the projection $A\times(C/c)\to A$ has a right adjoint and belongs to $W$, and likewise for $B$. Two-out-of-three makes the local comparison a member of $W$; relative Theorem A then gives $f\times\id_C\in W$. Factoring a finite product into maps changing one coordinate proves finite-product stability. This proof specifically uses the relative axiom at its last step.

The source's phrase that infinite stability holds for practically no localizer should be read as an informal warning, not a quantified theorem. There are immediate degenerate exceptions, such as the class of all functors. Ordinary weak equivalences do not in general enjoy the asserted strict-product stability for arbitrary families of categories: \citet[pp.~306--308]{Thomason1980} exhibits the discrepancy using products of a category with two parallel arrows. The lesson needed here is simply that the finite argument must not be extended to infinitely many coordinate changes.
\end{dcomment}

\begin{dcomment}{development}{p10}{A missing hypothesis in the product lemma}
\textbf{A finite counterexample to the literal generality of Lemma 2.5.}\ours
The qualification ``$W$ an arbitrary subset'' already causes a problem for binary products, independently of the editors' infinite-product warning. Let $\mathcal M$ be the four-element Boolean lattice
\[
0<a<1,\qquad 0<b<1,
\]
with $a$ and $b$ incomparable, regarded as a category, and put
\[
W=\{\id_1,\ a\longrightarrow1\}.
\]
All small products exist and are meets. A product of arrows in $W$ is $\id_1$ if all factors are identities, and is $a\to1$ otherwise; this includes the empty product. Thus even the strongest product-closure hypothesis printed in the lemma holds.

After localization $a$ and $1$ are isomorphic, but $0$ and $b$ are not. To verify the latter assertion without guessing the localized category, use the functor to the chain $0<b<1$ sending $a$ to $1$ and fixing $0,b,1$. It inverts $W$, hence factors through the localization, and distinguishes $0$ from $b$. Since $a\times b=0$ in $\mathcal M$, preservation of binary products would make the image of $0$ a product of the images of $1$ and $b$. The image of $1$ remains terminal, so this product is the image of $b$, a contradiction. This elementary counterexample is included to repair the wording, with no claim of historical priority.

The appropriate finite version assumes that $W$ contains all identities, as explicitly in \citet[Lemma 2.1.7 and Proposition 2.1.8, pp.~114--115]{Maltsiniotis2020}. For finite $I$, localization then commutes with the product of relative categories:
\[
(\mathcal M^I)[(W^I)^{-1}]
 \simeq (\mathcal M[W^{-1}])^I.
\]
One can prove this by changing one coordinate at a time. Every formal inverse in one coordinate lifts to an inverse of a tuple with identities in the other coordinates; coordinate changes commute, and finitely many finite words combine into one finite word. If products preserve $W$, the diagonal/product adjunction therefore descends and proves the corrected finite lemma. The same reasoning applies to coproducts. Every localizer under consideration in the chapter contains identities, so the counterexample does not affect those finite applications.

For infinite $I$, coordinatewise finite words need not have a common finite length; the displayed comparison is no longer a formal consequence of the localization universal property. We do not claim a general resolution of the infinite variant under just its printed assumptions. In a Quillen model category the editors' repair works: the product of fibrant replacements computes the derived product, and assumed closure of $W$ under these products identifies it with the ordinary product. Proposition 2.4 should accordingly be read with the editors' strong-localization hypothesis for finite products and with a separate justification whenever arbitrary coproducts are asserted; the finite adjunction argument alone is insufficient for an infinite family.
\end{dcomment}

\begin{dcomment}{exposition}{p11}{What descends from an adjunction}
Lemma 2.6 itself has no product-size defect. Suppose $F:\mathcal M\rightleftarrows\mathcal N:G$ is an adjunction and both functors send the designated equivalences into equivalences. Its unit and counit give natural transformations after localization. To check naturality on a formally inverted arrow, invert the corresponding naturality square; thus naturality on ordinary arrows suffices. The two triangle identities remain identities after applying the localization functors. They prove $\overline F\dashv\overline G$ directly, without a calculus of fractions or any claim that localization commutes with an infinite product of categories.

This also proves the final assertion of Lemma 2.5 for an arbitrary class $W$: a terminal object remains terminal. Use the adjunction between the unique functor $\mathcal M\to e$ and the selection of its terminal object, allowing the image of designated arrows to be merely inverted, as the universal property permits. Dually an initial object remains initial. Strictness is another matter; localization may create arrows into an initial object by reversing arrows out of it. The hypotheses of Lemma 2.14 are designed precisely to prevent that phenomenon. Keeping existence, strictness, and preservation of an entire family of products separate prevents several otherwise tempting shortcuts in the next pages.
\end{dcomment}

\begin{dcomment}{status}{p13:je renvoie à XV 9.12 (p. 133 bis).}{Connectivity and a hypothesis hidden in the recalled proof}
The equivalence of (2.7a) and (2.7b) is exactly the universal property of localization: $\pi_0$ factors if and only if it inverts $W$. The remaining equivalences require more care than the reference to XV might suggest. In the proof of XV, Proposition 9.12, ms.~133 quinquies, the editors explicitly flag the use of finite-product preservation as unproved under the weak axioms (1,2,3); it is available under (3 bis) \citep{Grothendieck1990XV}. Thus XVI, Proposition 2.7 inherits a missing hypothesis in the cited proof. We do not turn this gap into either a counterexample or an assertion that the weak formulation is an established theorem.

Here is the complete argument once finite-product preservation is supplied, in particular under (3 bis). Suppose $f:A\to B$ belongs to $W$ but does not induce a bijection on components. If two different components of $A$ are identified, choose points in them and a map $A\to\{0,1\}$ separating them. Their images in $B$ are connected, hence equal as points after localization, so invertibility of $f$ makes the two sections of $\{0,1\}\to e$ equal in $\Hot_W$. If a component of $B$ is missed, compare the constant map $B\to\{0,1\}$ with a map which is $1$ on that component and $0$ on the image of $A$. Their composites with $f$ agree. Inverting $f$ and evaluating at a point in the missed component gives the same equality of sections.

Product preservation now makes the two inclusions $X\to X\amalg X$ equal for every $X$. Composing with a map specified by any two functors $X\rightrightarrows Y$ shows that their localized images are equal. For nonempty $X$, a chosen point $e\to X$ is therefore inverse to $X\to e$. The localization is equivalent to $\Delta^1$ if no arrow from the empty category to a nonempty category is inverted, and to $e$ otherwise. Conversely, if $\pi_0$ factors, the discrete categories with different numbers of objects remain distinguishable, so neither degeneration occurs.

Corollary 2.8's forward implication has no such gap: $X\to e$ in $W$ immediately gives $\pi_0X=\{*\}$. Its converse, as argued through the recalled classification and strong saturation, needs the same repair to that classification. The test of Corollary 2.9 has a direct proof independent of it. Pull a $W$-map back to each component of its target using (4b); the stated connected-target condition says that every inverse image is nonempty and connected, exactly the bijectivity of $\pi_0$. The rest of this section explicitly assumes (5), so the concrete description of points below does not depend on the unresolved weak-hypothesis step.
\end{dcomment}

\begin{dcomment}{exposition}{p15:La formule d'adjonction (2.8.2) est d'ailleurs familière}{The double adjunction without an infinite-coproduct shortcut}
Write $d(I)=I_W$ for the image of the discrete category on a set $I$. The first adjunction $\pi_0\dashv d$ follows by localizing the familiar adjunction on $\Cat$: every functor to a discrete category is constant on components. The second,
\[
d\dashv\pi_0,\qquad
\Hom_{\Hot_W}(dI,X)=\operatorname{Set}(I,\pi_0X),
\]
is special to the localized situation. In the printed proof it is reduced to $I=e$ by treating $dI$ as the coproduct of $I$ copies of $e$. For arbitrary $I$ that reduction calls on the earlier infinite-coproduct assertion. The following direct proof avoids it and works with just (1,2,3,5).

First suppose two functors $u,v:I\to X$ choose points in the same component for every $i$. For each $i$ choose a finite zigzag joining $u(i)$ to $v(i)$, and let $J_i$ be its walking zigzag category. Its nerve is contractible, so $J_i\to e$ is in $W_\infty\subset W$; alternatively this follows by successive elementary contractions. Put $J=\coprod_i J_i$. The projection $J\to I$ has comma category $J_i$ over $i$, and (3) therefore puts it in $W$ directly, without (4). Its two endpoint sections are both inverse to it in the localization. Composing with the chosen functor $J\to X$ proves $L(u)=L(v)$. The lengths of the individual zigzags need not be bounded.

Next let $w:X\to Y$ belong to $W$ and let $u:I\to Y$. Surjectivity of $\pi_0(w)$ allows a choice $s:I\to X$ such that $ws$ and $u$ lie in the same component pointwise. The previous paragraph gives $L(w)L(s)=L(u)$. Any localization morphism from $dI$ is represented by a finite word of ordinary arrows and inverses of arrows in $W$. Whenever its first reversed arrow occurs, this lifting replaces the preceding part of the word by an ordinary map from $I$. Induction eliminates all reversed arrows. Thus every morphism $dI\to X$ is represented by an actual functor $I\to X$; two representatives agree exactly when their component functions agree, since the descended $\pi_0$ detects that distinction. This proves the desired natural bijection, including $I=\emptyset$.

Taking $I=e$ proves Corollaries 2.9 and 2.12. Taking $X=dJ$ shows that $d$ is fully faithful. It also proves, rather than assumes, that $dI$ is the coproduct of $I$ copies of $e$ in $\Hot_W$, since its maps to $X$ are the $I$-indexed families of points of $X$. Here $I\times e$ denotes a \emph{copower}, as the source says; it is not an assertion that an external set is already an object participating in a categorical product.
\end{dcomment}

\begin{dcomment}{examples}{p17:En effet, cette bijectivité est évidente quand on travaille dans Cat,}{Why the two adjunctions are exceptional}
For a category $\mathcal H$ with terminal object $e$ and the required coproducts, the copower functor $dI=\coprod_I e$ is automatically left adjoint to $\Gamma=\Hom(e,-)$. The unit $I\to\Gamma(dI)$ is bijective precisely when $d$ is fully faithful. Neither this condition nor the existence of coproducts automatically supplies the other adjunction: that further assertion is that every function $\Gamma(X)\to I$ comes from exactly one morphism $X\to dI$.

The examples distinguish the conditions. In sets both adjunctions hold with $d=\Gamma=\id$. In $\Cat$, $d$ is fully faithful and $\pi_0\dashv d$, but $\Gamma(X)=\Ob(X)$ is usually different from $\pi_0X$; two objects joined by an arrow still define different functors $e\to X$. Localization makes those two functors equal. In an additive category the terminal object is the zero object, so every copower of it is again zero and $\Gamma(X)$ is a singleton. Consequently $I\to\Gamma(dI)$ is bijective only for singleton $I$. This explains the warning in footnote 6 without treating additivity as a mild exceptional case of the asserted unit condition.

In the triangular display the edition prints the horizontal top arrow as well as the two dotted downward arrows, the right one labelled $\sim$; the transcription keeps all three. The top arrow is the composite of the left arrow with the inverse of the right arrow, exactly as the surrounding prose requires.
\end{dcomment}

\begin{dcomment}{exposition}{p15'}{Shortening localization words and preserving emptiness}
Lemma 2.14 is a useful statement about localization words, rather than a calculus-of-fractions theorem. A morphism is a finite zigzag; it need not have a single denominator. If a zigzag ends at a strict initial object $0$, the prohibition on $W$-arrows $0\to X$ with $X\ne0$ forces the adjacent object to be $0$. Strictness then forces the preceding object to be $0$, and induction runs through the zigzag. This proves strictness after localization. The repeated subscript $2n-1$ in the printed proof is an edition slip: the second conclusion is $X_{2n-2}=0$.

For a word beginning at $e$, the lifting condition replaces $e\to X_1\leftarrow X_2$ by a genuine map $e\to X_2$ having the same localized composite. Composing with the next arrow shortens the word by two. In $\Cat$, (5) supplies this lifting because every component of $X_1$ has a representative in $X_2$, while a zigzag inside a component becomes equality of points after localization. The citation of Corollary 2.12 when checking condition (1) is circular as printed; that condition follows directly from (5), or from the empty-target argument of Corollary 2.13. Neither slip affects the corrected induction.

There is a useful logical division in the final commentary. The class of all functors satisfies every purely closure-type axiom under discussion, so no conjunction of those axioms can force (5). Conversely, sum closure need not be an independent axiom once a gluing theory is available: a coproduct is the union of two disjoint open subcategories, giving finite closure from (6a), and its finite partial unions give arbitrary closure after adding filtered-colimit stability. Reflection on a summand still comes from the separate base-change requirement. This is why (6a,7,8), rather than just (6a,8), appears in the correct final implication to (4).
\end{dcomment}

\begin{dcomment}{exposition}{p16'}{The three cube rules measure different phenomena}
An open subcategory here is a sieve: whenever an arrow has target in it, its source belongs to it as well. If $X=U\cup V$ is a cover by two sieves and $S=U\cap V$, every nerve simplex lies wholly in $U$ or wholly in $V$, by looking at its last vertex. Therefore
\[
NX=NU\amalg_{NS}NV.
\]
Both inclusions are simplicial monomorphisms, so this strict simplicial pushout also represents the homotopy pushout of spaces. For cosieves the same reasoning uses the first vertex. This elementary geometry underlies the forward cube rule (6a).

The reverse rule (6b) is a different assertion: equivalences on the two pieces and the union must detect an equivalence on the intersection. A homotopy pushout does not ordinarily have this property. The third rule (6c) asks for detection on both pieces from the intersection and union. It is consequently misleading to call all three rules merely different forms of the ordinary gluing lemma. Their implications require arguments and, for the converse involving (6c), the stated retract hypothesis.

The manuscript's evolving terminology should be followed literally. Footnotes 11 and 13 remove (6c) from the name Loc (6), but Loc (6) still means \emph{both} (6a) and (6b). Only (6a) is the familiar forward gluing property of ordinary weak homotopy equivalences. The commentary to XIIb-a, ms.~78--79, and XII Appendix, ms.~6--17, already separates these rules and gives the acyclic-link obstruction to reverse excision for $W_\infty$; those examples need not be duplicated here. The finite-dimensional examples below settle the particular prediction made on the present page.

The small cubes in Proposition 3.1 can be checked by recording their vertical tuples. For $Q$ the tuple is $(\id,i_1,\id,j_1)$; for $Q'$ it is $(i_1,\id,j_1,\id)$. Substitution immediately gives all three implications. The printed second occurrence of $f_1$ in the explanation of $Q'$ means $f_2$. The doubtful $\id$ label belongs to the front vertical edge $X_1\to X_1$, as the diagram and the edition both confirm.
\end{dcomment}

\begin{dcomment}{examples}{p17'}{A reverse-excision counterexample in every finite degree}
Let $W_n$ consist of functors inducing a bijection on components and isomorphisms on homotopy groups through degree $n$, for every choice of base point. Grothendieck's prediction is correct for every finite $n\geq0$: $W_n$ satisfies forward gluing, but fails (6b).

Here is one uniform finite example. Let $K$ be the boundary of an $(n+1)$-simplex, so $|K|=S^n$, and take its two simplicial cones $a*K$ and $b*K$. For a simplicial complex $L$, let $F(L)$ be the poset of its nonempty faces, ordered by inclusion. Its nerve is the barycentric subdivision of $L$. The inclusions
\[
F(K)\subset F(a*K),\qquad F(K)\subset F(b*K)
\]
are sieves, and the two cone face posets cover
\[
F((a*K)\cup_K(b*K)).
\]
Thus they form an MV square. Map this entire square to the square all of whose vertices are $e$. The two piece maps are ordinary weak equivalences because cones are contractible. The union has homotopy type $S^{n+1}$, so its map to $e$ belongs to $W_n$. The intersection has homotopy type $S^n$, so its map to $e$ does not belong to $W_n$: for $n=0$ it has two components, and for $n>0$ its $n$th homotopy group is $\mathbb Z$. This is exactly a failure of (6b), using finite posets throughout.

For completeness, the forward assertion can be seen from the universal property of $n$-truncation. A map is an $n$-equivalence exactly when it induces equivalences of derived mapping spaces into all $n$-types. Mapping from a homotopy pushout into an $n$-type gives a homotopy pullback of mapping spaces. Equivalences on its three entries therefore give an equivalence on the result. Since an MV square has the homotopy-pushout realization just proved, (6a) follows. This proves the source's conjectural example rather than merely restating it. Notice that passage from $S^n$ to its suspension is the mechanism of failure: the union moves the first detectable homotopy group one degree beyond the range being tested.
\end{dcomment}

\begin{dcomment}{exposition}{p19}{The Dwyer neighborhood and the missing gluing step}
A Dwyer inclusion $A\hookrightarrow B$ is a sieve whose inclusion into its generated cosieve $U$ has a right adjoint $r:U\to A$ restricting to the identity on $A$. Thus $A\to U$ is a categorical homotopy equivalence. Put $V=B\setminus A$ and $T=U\cap V$. Both $U$ and $V$ are cosieves, and $B=U\cup_T V$. On pushing $A\to B$ out along $A\to A'$, one obtains the corresponding cover
\[
B'=U'\cup_T V,\qquad U'=U\amalg_A A'.
\]
The inclusion $A'\to U'$ again has a right-adjoint retraction, induced by $r$, so is in every $W$. This is the neighborhood calculation in \citet[Definition 4.1 and Proposition 4.3, pp.~313--317]{Thomason1980}; its proof does not require the claim elsewhere in that paper that all retracts of Dwyer maps are Dwyer maps, disproved in \citet{Cisinski1999Dwyer}.

It explains both uses of the excision axioms here. If $A\to A'$ is in $W$, then $U\to U'$ is in $W$ by the two homotopy-equivalent neighborhoods. The cube from the first closed cover to the second has its intersection and $V$ maps equal to identities. The restricted case of (6a) therefore gives $B\to B'\in W$. This proves that Dwyer inclusions are intrinsic cofibrations, since their pushouts are again Dwyer inclusions.

Alternatively, suppose $A\to B$ itself is in $W$ and assume (6$'$a,b). Then $U\to B$ is in $W$. The reverse single-square rule, applied to $B=U\cup_T V$, gives $T\to V\in W$. The forward rule in the new cover gives $U'\to B'\in W$, hence $A'\to B'\in W$. Thus trivial Dwyer inclusions are stable under every pushout. Opposite invariance justifies using closed covers in both arguments. This is precisely condition (4) in Brown's lemma, and makes explicit where both directions of (6$'$) enter the second proof.

The other factorization hypothesis also has a concrete explanation. Let $J=(0\to1\leftarrow2)$ and form
\[
M(f)=(A\times J)\amalg_{A\times\{2\}}B
\]
for $f:A\to B$. The inclusion of $A\times\{0\}$ is Dwyer: its cosieve neighborhood is $A\times(0\to1)$, with right-adjoint projection to $A$. The projection $M(f)\to B$ is a categorical homotopy equivalence. Indeed, the endomorphism of $J$ sending $0$ to $1$ and fixing $1,2$ receives natural transformations from both $\id_J$ and the constant functor at $2$, stationary at $2$; these descend after the pushout and contract $M(f)$ onto $B$. Hence every functor factors as a Dwyer inclusion followed by a member of $W$. No transfinite construction or functorial factorization axiom is needed for this particular factorization.
\end{dcomment}

\begin{dcomment}{exposition}{p21}{A complete proof of the Brown lemma used here}
The relevant historical result is the dual of the properness argument for categories of fibrant objects, not just the isolated factorization lemma: see \citet[Part I, \S1, pp.~420--422, and \S4, Lemmas 1--2, pp.~428--429]{Brown1973}. A modern one-sided formulation appears in \citet[\S\S1.3--1.4]{RadulescuBanu2009}. Here is a proof directly under the four displayed hypotheses, filling the gap Grothendieck explicitly acknowledges on ms.~22.

\emph{First, Brown's functor criterion.} A functor $F$ from such a category to a category with a two-out-of-three class of equivalences preserves all $W$-maps if it preserves $C\cap W$. For $w:X\to Y$ in $W$, factor
\[
(w,\id_Y):X\amalg Y\xrightarrow{k} Z\xrightarrow{p}Y,
\qquad k\in C,\ p\in W.
\]
The restrictions $k_X,k_Y$ belong to $C$, because coproduct inclusions do, and both are in $W$ by two-out-of-three. Their images under $F$ are equivalences. Since $pk_Y=\id_Y$, $F(p)$ is an equivalence, and so is $F(w)=F(p)F(k_X)$.

\emph{Second, change the initial object.} For each $A$, consider the category $\mathcal C_A$ of arrows $A\to X$ belonging to $C$, with all maps under $A$. Take equivalences and cofibrations to be those whose underlying maps are in $W$ and $C$. It satisfies the same hypotheses: its initial object is $\id_A$, pushouts and factorizations are inherited, and composing a factorization's first map with $A\to X$ keeps the new object in $\mathcal C_A$. For any $u:A\to A'$, the functor
\[
u_!: \mathcal C_A\longrightarrow\mathcal C_{A'},
\qquad (A\to X)\longmapsto(A'\to A'\amalg_A X)
\]
preserves trivial cofibrations by condition (4). The functor criterion shows that it preserves all equivalences. This assertion concerns maps \emph{between} cofibrant objects under $A$, and is not yet the desired properness statement; the distinction prevents a circular proof.

\emph{Third, prove properness.} We must show that the pushout of $w:A\to A'$ in $W$ along $i:A\to B$ in $C$ is a $W$-map $B\to P$. Applying the factorization of the first paragraph to $(w,\id_{A'})$ writes $w=pj$, where $j\in C\cap W$ and $p$ has a section $s\in C\cap W$. Pushing out $j$ is covered by (4). We are reduced to the case $ws=\id_{A'}$, with $s:A'\to A$ a trivial cofibration.

Form $B_1=A\amalg_{A'}B$, using $s:A'\to A$ and $is:A'\to B$. The map $q:B\to B_1$ is a pushout of $s$, hence a trivial cofibration. There is a retraction $g:B_1\to B$ induced by $(i,\id_B)$, so $g\in W$. The structural maps $A\to B_1$ and $A\to B$ belong to $C$: the former is a pushout of $is\in C$. Consequently $g$ is an equivalence in $\mathcal C_A$. Apply $w_!$. Since $ws=\id_{A'}$, associativity of pushouts gives
\[
A'\amalg_A B_1\simeq B,
\qquad A'\amalg_A B=P.
\]
Under these identifications, $w_!(g)$ is precisely $B\to P$, which is therefore in $W$. Stability of $C$ under pushouts makes the same argument apply to every cobase change of $i$, proving $C\subset\Cof_W$.

For the converse in Corollary 3.2.2, suppose $C\subset\Cof_W$, let $i:A\to B$ be in $C\cap W$, and push it out along an arbitrary $f:A\to A'$. Factor $f=pj$ with $j\in C$ and $p\in W$. Pushing $i$ out along $j$ gives a $W$-map because $j$ is an intrinsic cofibration. Pushing $p$ out along that new $C$-map gives another $W$-map. Two-out-of-three, applied to their composite and $p$, proves that the final pushout of $i$ is in $W$. This verifies all four equivalent conditions without assuming the equality $\Cof_W\cap W=W^{\univ}$ as an unexplained extra theorem. The argument only uses identities and two-out-of-three in $W$; retract closure and functorial factorizations are unnecessary here.
\end{dcomment}

\begin{dcomment}{status}{p25}{What the excision reformulations establish}
The new contribution of the present argument, relative to the earlier cube calculations in XIIb-a and XII Appendix, is the missing cofibration criterion. The Dwyer-neighborhood proof and Brown's lemma establish it from (6$'$a,b), so the earlier passage from single-square reversibility to the two cube rules can now be used without assuming the desired cofibration theory in advance. The source's equivalence
\[
(6\mathrm{a},6\mathrm{b})\Longleftrightarrow(6'\mathrm{a},6'\mathrm{b})
\]
should be read with that repair. Closure of $W$ under retracts is a separate premise for the converse implication to (6c); it concerns $W$, not closure of the class of Dwyer inclusions under retracts. The latter closure is false \citep[pp.~228--229]{Cisinski1999Dwyer}.

There is also a short direct check of the cover comparison in question (1$^\circ$ bis). Let $G$ be the Grothendieck construction on the span $S\to U$, $S\to V$, and $q:G\to X$ its canonical functor. For $x\in U$, the sieve condition gives $U/x=X/x$ and $V/x=S/x$. Thus $G/x$ is the integral of the span $S/x\to X/x$, $S/x\xrightarrow{\id}S/x$. The redundant branch contracts by a zigzag of natural transformations, leaving $X/x$, which has terminal object $\id_x$. The same argument applies if $x\in V$. Hence every $G/x$ is contractible and $q\in W$ by (3). This proves the cover comparison without (3 bis), in agreement with the revised proof mentioned in footnote 17. For general spans, identifying their integral with the intrinsic homotopy pushout still uses the cofibration theory; the earlier commentary to XIIb-b, ms.~120--149, discusses that distinction in detail.

The statements about monomorphisms of presheaves on ms.~24 refer to $W_A$ as defined using the category of elements. They are claims about the induced relative category, not a claim that every fully faithful functor in $\Cat$ is a cofibration. Nor does ``enough intrinsic cofibrations'' by itself provide both lifting systems of a Quillen model structure. The conclusion licensed at this stage is Brown's one-sided cofibration theory, with its homotopy pushouts. Proposition 3.6 and footnote 19 further reduce the issue to enough such cofibrations; the promised implication (6A)$\Rightarrow$(6B) is addressed at ms.~66--68, not proved at its first announcement here.

The bare implication (6$'$a)$\Rightarrow$(6a) is explicitly left as a question in this passage; footnote 18 supplies the additional assumptions (8,9), not an unconditional answer. Likewise the final suggestion that (6c) might force retract closure is a conjectural observation. For strong localizers retract closure follows independently from later saturation theory, but that does not answer the question for the weak framework of Definition 1.1. No result settling these precise weakest-hypothesis formulations was found in the sources checked. The next two blocks instead extract a definite, and rather restrictive, consequence of the reverse cube axiom itself.
\end{dcomment}

\begin{dcomment}{development}{p27:je soupçonne que Loc (6 c)}{Integral homology supplies a reversible localizer}
\textbf{Homological Theorem A and reversible excision.}\ours
Let $C_*(A)=C_*(NA;\mathbb Z)$ be the integral nerve chain complex and define
\[
W_{H\mathbb Z}=\{f:A\to B\mid H_k(NA;\mathbb Z)\xrightarrow{\sim}
H_k(NB;\mathbb Z)\text{ for all }k\geq0\}.
\]
This is a strongly saturated fundamental localizer satisfying (3 bis), (5), and all three cube rules (6a,b,c). We include the proof because the ensuing minimality statement needs the actual Theorem A hypothesis, not just the observation that homology has a Mayer--Vietoris sequence. The ingredients are classical; the formulation and proof assembled here are offered as a development of this passage, with no priority claim.

\emph{The local-to-global assertion.} For $u:A\to S$, use the first-quadrant double complex $D(u)_{p,q}$ freely generated by strings
\[
a_0\longrightarrow\cdots\longrightarrow a_q,
\qquad u(a_q)\longrightarrow s_0\longrightarrow\cdots\longrightarrow s_p.
\]
The vertical faces compose or omit arrows in the $A$-string, including composition with the connecting arrow when $a_q$ is omitted. The horizontal faces do the same for the $S$-string. Use unnormalized chains for this description and total differential $d_v+(-1)^q d_h$. There is an augmentation
\[
r_A:\operatorname{Tot}D(u)\longrightarrow C_*(A)
\]
which forgets $s_0$ in horizontal degree zero and is zero in positive horizontal degrees. For a fixed $A$-string, the horizontal complex is that of $u(a_q)\backslash S$. Its initial object is $\id_{u(a_q)}$, so its augmented chain complex is contractible. Filtering the total complex therefore proves that $r_A$ is a quasi-isomorphism for every $u$.

Now take a triangle $A\xrightarrow{f}B\to S$. It induces $D(A\to S)\to D(B\to S)$. For a fixed string $s_0\to\cdots\to s_p$, the corresponding vertical complexes are precisely $C_*(A/s_0)$ and $C_*(B/s_0)$. If all the comma comparisons are integral homology equivalences, this map is a quasi-isomorphism in every such column. The other first-quadrant filtration consequently makes the map on totals a quasi-isomorphism. There is no convergence issue: in each total degree only finitely many bidegrees occur. The naturality square for the two augmentations $r_A,r_B$ identifies the resulting map with $C_*(f)$, proving (3 bis). Taking $B=S$ and using the terminal objects of $S/s$ also proves the weak axiom (3).

For readers wishing to check the comparison with the ordinary homological Theorem A directly, there is a second augmentation $r_S:\operatorname{Tot}D(u)\to C_*(S)$, forgetting the $A$-string in vertical degree zero. The signed concatenation
\[
h(a_\bullet,u(a_q)\to s_\bullet)
 =(-1)^q\bigl[u(a_0)\to\cdots\to u(a_q)\to s_0\to\cdots\to s_p\bigr]
\]
satisfies $dh+hd=r_S-u_*r_A$. The internal faces cancel in pairs; deleting the sole $A$-vertex when $q=0$ yields $r_S$, and deleting the sole $S$-vertex when $p=0$ yields $-u_*r_A$. If each $A/s$ has the homology of a point, the vertical filtration makes $r_S$ a quasi-isomorphism as well. This explicit homotopy verifies that the isomorphism produced by the two filtrations is induced by the given functor, rather than merely an abstract isomorphism between homology groups.

\emph{The other localizer axioms.} A terminal object contracts the nerve, giving (2). Identities, two-out-of-three, mild saturation, and retract closure follow by applying integral homology. Strong saturation follows too: homology factors through $\Cat[W_{H\mathbb Z}^{-1}]$, so any arrow inverted in that localization already induces a homology isomorphism. In degree zero, a functor induces the free abelian map on its function of components; such a map is an isomorphism exactly when the function is a bijection. This proves (5).

\emph{The three cube rules.} For an open MV square, the nerve decomposition gives a natural short exact sequence of normalized chain complexes
\[
0\longrightarrow C_*(S)
 \xrightarrow{(i_*,-j_*)}C_*(U)\oplus C_*(V)
 \longrightarrow C_*(X)\longrightarrow0.
\]
A map of MV squares gives a map of these short exact sequences. If two of the three vertical maps are quasi-isomorphisms, so is the third, by their long exact homology sequences. This yields (6a) and (6b). For (6c), the two outer quasi-isomorphisms force the middle map to be a quasi-isomorphism. That middle map is block diagonal, so each of its two summand maps is a quasi-isomorphism separately. These arguments apply to disconnected and empty categories as written. The same proof works for closed covers, or one may pass to opposite categories. Thus the homological class is an actual example of the reversible theory sought here, rather than an analogy with an exact sequence.
\end{dcomment}

\begin{dcomment}{development}{p27:je soupçonne que Loc (6 c)}{The least localizer satisfying reverse cube excision}
\textbf{Minimality forced by one suspension.}\ours
Among weak fundamental localizers in the sense of Definition 1.1 satisfying (6b), the least one is $W_{H\mathbb Z}$. In particular, the least such localizer satisfying both directions of Loc (6), or satisfying all three cube rules, is again $W_{H\mathbb Z}$. The previous block proves the required upper bound. We now prove that every such $W$ contains every integral homology equivalence. This gives a precise consequence of Grothendieck's reversibility requirement; it does not claim that every larger reversible weak localizer is strong or saturated.

By \citet[Theorem 6.1.18]{Cisinski2006}, $W_\infty\subset W$ even under the weak axioms. Let $f:A\to B$ be an integral homology equivalence. Its map on components is a bijection, and its restriction to each matched pair of components is a homology equivalence: homology is the direct sum over components and the induced map respects this decomposition. Write these maps as $f_i:A_i\to B_i$, with each category nonempty and connected. If there are no components, $f$ is $\id_\emptyset$ and the assertion is immediate.

For a nonempty category $C$, form $C_-$ and $C_+$ by adjoining one new terminal object in each case, and let $\Sigma C$ be the category obtained by adjoining both objects, with no morphisms between them. Then
\[
\begin{tikzcd}[column sep=large]
C\arrow[r]\arrow[d] & C_-\arrow[d]\\
C_+\arrow[r] & \Sigma C
\end{tikzcd}
\]
is an open MV square: both cones are sieves, their intersection is $C$, and their union is $\Sigma C$. The nerve realizes the union of the two ordinary cones on $|NC|$, hence the unreduced suspension of $|NC|$. Each cone is contractible. Because $C$ is connected, its suspension is simply connected, by van Kampen, and the suspension isomorphism identifies its reduced integral homology with that of $C$, shifted one degree.

Apply this construction to each $f_i$. The map on either cone is in $W_\infty$. The map $|N\Sigma A_i|\to|N\Sigma B_i|$ is a homology equivalence between simply connected CW complexes. The homological Whitehead theorem makes it a weak homotopy equivalence \citep[Corollary 4.33, pp.~367--368]{Hatcher2002}. Thus the three maps on the pieces and union belong to $W_\infty\subset W$. Reverse cube excision (6b) forces $f_i\in W$.

One must not finish by taking a sum of these conclusions unless sum closure has been assumed. Instead take the disjoint union of the entire componentwise cube diagrams \emph{before} applying (6b). This is a single open MV cube whose intersection map is the original $f$. The two cone maps and the union map are ordinary weak equivalences globally: their component functions are the same bijection, and their restrictions are the weak equivalences already established. Therefore they belong to $W$ without any use of (4), and (6b) directly gives $f\in W$. This proves the lower bound and the theorem.

The conclusion also locates the obstruction to identifying reversible excision with ordinary homotopy theory. Noncontractible acyclic spaces, such as the example recalled in \citet[discussion following Corollary 4.33, p.~368]{Hatcher2002}, have the homology of a point but a nontrivial fundamental group; categorical models give maps in $W_{H\mathbb Z}\setminus W_\infty$. Such examples were already used in the commentary on Chapter~XII to disprove reverse excision for $W_\infty$. The new content here is the universal lower bound: reverse cube excision forces \emph{every} integral homology equivalence, not just the particular acyclic examples. We did not find this precise least-localizer formulation in the primary sources examined; its proof is included in full so that the assertion can be assessed independently of any priority question.

\textbf{Proposition (the single-square form and the cofibre test).}\ours
\textit{Let $W$ be a weak fundamental localizer satisfying the reverse implication for one cover by two sieves, in the form of W(7 bis) of Chapter~XII, ms.~76: if $X=U\cup V$ with $U,V$ sieves and $U\to X$ is in $W$, then $U\cap V\to V$ is in $W$. For a functor $g:A\to B$ let $M_g$ be the category made of $A$, of $B$, and of the arrows $a\to b$ given by the arrows $g(a)\to b$ of $B$, and let $K(g)$ be $M_g$ with one more object $t$, one arrow $a\to t$ for every object $a$ of $A$, and no other new arrow than $\id_t$. Then: (a) if $K(g)$ is $W$-aspherical, $g\in W$; (b) $W\supset W_{H\mathbb Z}$.}

\textit{Proof.}
The subcategories $M_g$ and $A_+=A\cup\{t\}$ of $K(g)$ are sieves, they cover $K(g)$, and they meet in $A$. The category $A_+$ has the terminal object $t$. If $K(g)$ is $W$-aspherical, $A_+\to K(g)$ is in $W$ by two-out-of-three, and the hypothesis gives $A\to M_g$ in $W$. The retraction $M_g\to B$ which is $g$ on $A$ has the inclusion of $B$ as right adjoint, so it is in $W_\infty\subset W$; its composite with $A\to M_g$ is $g$. This proves (a).

For (b) let $g$ be an integral homology equivalence. The nerve of $A_+$ is the cone on $NA$, so $BK(g)$ is the mapping cone of $Bg$. The function $\pi_0(g)$ is bijective; hence $C=K(g)$ is connected (it is the point $t$ if $A$ and $B$ are empty), and it is acyclic by the exact sequence of the cone. The space $BK(C\to e)$ is the unreduced suspension of $BC$: simply connected and acyclic, hence weakly contractible. So $K(C\to e)$ is $W_\infty$-aspherical, hence $W$-aspherical. By (a) for $C\to e$, the category $C$ is $W$-aspherical; by (a) for $g$, $g\in W$.\hfill$\square$

The proposition replaces the cube axiom (6b) of the previous argument by its single-square form. For localizers satisfying the relative axiom (3 bis), forward gluing gives the converse of (a), so that $g\in W$ if and only if $K(g)$ is $W$-aspherical, and this criterion is in turn equivalent to the cube forms (6b) and (6c); that equivalence is not reproduced here.

\textbf{The reversible localizers are not the localizers containing $W_{H\mathbb Z}$.}\ours
The localizer $W_0$ of the functors bijective on connected components contains $W_{H\mathbb Z}$ and does not satisfy the reverse implication. For the fold $g:\{0,1\}\to e$ the category $K(g)$ has the four objects $0,1,e,t$ and the arrows $0\to e$, $1\to e$, $0\to t$, $1\to t$; it is connected, so $W_0$-aspherical, while $g\notin W_0$. By (a), $W_0$ violates W(7 bis). The class of localizers with reverse excision has the least element $W_{H\mathbb Z}$ and is not closed upwards.
\end{dcomment}

\begin{dcomment}{development}{p27:je soupçonne que Loc (6 c)}{Reverse excision excludes the fibration axioms}
Reverse excision and the fibration axioms of Sections 4--6 pull in opposite directions. In this block a \emph{reversible localizer} is a fundamental localizer satisfying the relative axiom (3 bis) and the reverse implication for one cover by two sieves, as in the proposition just proved; part (a) of that proposition is called the cofibre test. Such a localizer is strongly saturated \citep[Proposition~4.2.4]{Cisinski2006}, and it contains $W_{H\mathbb Z}$. The statements below were not found in the sources consulted; no priority is claimed.

\textbf{Proposition 1 (no sphere is aspherical).}\ours
\textit{Let $W$ be a reversible localizer. If some category of the homotopy type of a sphere $S^n$, $n\geq0$, is $W$-aspherical, then $W=\operatorname{Fl}(\mathrm{Cat})$. In particular a reversible localizer other than $\operatorname{Fl}(\mathrm{Cat})$ does not contain the localizer $W_1$ of the functors inducing an equivalence of fundamental groupoids.}

\textit{Proof.}
Since $W\supset W_\infty$, asphericity depends only on the homotopy type. If $X$ has the type of $S^{n-1}$, $n\geq1$, then $K(X\to e)$ has the type of $S^n$; if it is aspherical, the cofibre test makes $X$ aspherical. Descending, the discrete category with two objects is $W$-aspherical. It is $K(\emptyset\to e)$, so $\emptyset\to e$ is in $W$ by the cofibre test. For a category $X$, the functor $\emptyset\to X$ is then in $W$ by (3 bis) over $X$: over each object $x$ it is $\emptyset\to X/x$, and $X/x$ is aspherical. Two-out-of-three gives $X\to e\in W$ for every $X$, and then $f\in W$ for every functor $f$. A category of the type of $S^2$ is $W_1$-aspherical; this gives the last assertion.\hfill$\square$

\textbf{Lemma (aspherical objects do not see free groups).}\ours
\textit{Let $W\neq\operatorname{Fl}(\mathrm{Cat})$ be a strongly saturated fundamental localizer containing $W_\infty$ for which no category of the type of $S^0$ or $S^1$ is aspherical, and let $F$ be a free group. For every $W$-aspherical category $A$ the space $BA$ is nonempty and connected, every homomorphism $\pi_1(BA)\to F$ is trivial, and the inclusion of the constant maps $BF\to\operatorname{Map}(BA,BF)$ is a weak equivalence.}

\textit{Proof.}
If $A$ were empty, $\emptyset\to e$ would be in $W$ and $W$ would be everything, as in the preceding proof. If $A$ were not connected, a surjection of $A$ onto the discrete category $\{0,1\}$ would have a section, the arrow $\{0,1\}\to e$ would be a retract of $A\to e$, and $S^0$ would be aspherical. Let $\varphi:\pi_1(BA)\to F$ be a nontrivial homomorphism. Its image is a nontrivial free group and maps onto $\mathbf Z$; a loop of $BA$ mapping to a generator and the map $BA\to S^1$ classifying the composite $\pi_1(BA)\to\mathbf Z$ make $S^1$ a retract of $BA$ in $\mathrm{Hot}$. In the localization by $W$ the image of $S^1\to e$ is then a retract of an isomorphism, hence an isomorphism, and strong saturation makes $S^1$ aspherical. So $\varphi$ is trivial. The components of $\operatorname{Map}(BA,BF)$ correspond to the conjugacy classes of homomorphisms $\pi_1(BA)\to F$, and the component of the trivial one is a $K(F,1)$ onto which the constants map by an equivalence.\hfill$\square$

\textbf{Proposition 2 (reverse excision against enough fibrations).}\ours
\textit{Let $W$ be a fundamental localizer satisfying (3 bis) which has enough intrinsic fibrations, Loc (7 F), and whose aspherical objects satisfy the conclusion of the lemma for the free group $F$ on two generators $x,y$. Then the functor $B\alpha:BF\to BF$ defined by the endomorphism $\alpha(x)=x$, $\alpha(y)=yxyx^{-1}y^{-1}$ is not in $W$. Consequently a reversible localizer with enough intrinsic fibrations is $\operatorname{Fl}(\mathrm{Cat})$.}

\textit{Proof.}
Let $W'$ be the class of functors $f$ for which $\operatorname{Map}(Bf,BF)$ is a weak equivalence. It is weakly saturated and contains $W_\infty$; it satisfies the relative Theorem~A, because $BA$ is the homotopy colimit of the spaces $B(A/c)$ over the base and $\operatorname{Map}(-,BF)$ turns homotopy colimits into homotopy limits. By hypothesis $A\to e$ is in $W'$ for every $W$-aspherical $A$.

We show $W\subset W'$. The class $\mathrm{Fib}_W$ is stable under composition and base change, every $X\to e$ is in it because $W$ is stable under products, and by (7 F) every functor is a $W$-equivalence followed by a member of $\mathrm{Fib}_W$. Brown's factorization lemma (commentary at ms.~21) then writes $f\in W$ as $f=ps$ with $p\in\mathrm{Fib}_W$ and $s$ a section of some $q\in\mathrm{Fib}_W\cap W$; so $p\in W$ too. A member $p:X\to Y$ of $\mathrm{Fib}_W\cap W$ stays in $W$ after every base change $Y'\to Y$: factor $Y'\to Y$ as $j:Y'\to Z$ in $W$ followed by $r:Z\to Y$ in $\mathrm{Fib}_W$; the pullback $p_Z$ of $p$ along $r$ is in $W$ because $r$ is an intrinsic fibration and $p\in W$; the map $X\times_YY'\to X\times_YZ$ is in $W$ because $p$ is an intrinsic fibration and $j\in W$; two-out-of-three gives the pullback of $p$ to $Y'$ in $W$. In particular $X/y\to Y/y$ is in $W$ for every object $y$, so $X/y$ is $W$-aspherical, so $X/y\to e$ is in $W'$, and Theorem~A for $W'$ gives $p\in W'$. The same holds for $q$; then $s\in W'$ and $f\in W'$.

The functor $B\alpha$ is not in $W'$. The components of $\operatorname{Map}(BF,BF)$ are the conjugacy classes of endomorphisms of $F$, and $B\alpha$ acts on them by $[\varphi]\mapsto[\varphi\alpha]$. If this were bijective there would be $\beta$ with $\beta\alpha$ inner; then $[\alpha\beta\alpha]=[\alpha]$, so $[\alpha\beta]=[\id]$ by injectivity, and $\alpha$ would be surjective. It is not: the homomorphisms $F\to S_3$ with $x\mapsto(2\,3)$ and $y\mapsto(2\,3)$, respectively $y\mapsto(1\,2)$, agree on $x$ and on $\alpha(y)$ (both send it to $(2\,3)$, since $(1\,2)(2\,3)$ has order three) and differ on $y$.

For the last assertion let $W\ne\operatorname{Fl}(\mathrm{Cat})$ be reversible. By Proposition~1 no sphere is $W$-aspherical, so the lemma applies. The endomorphism $\alpha$ is the identity on $F^{\mathrm{ab}}$, so $B\alpha$ is an integral homology equivalence and lies in $W_{H\mathbb Z}\subset W$. By the first part $W$ cannot satisfy (7 F).\hfill$\square$

\textbf{Consequences.}
(1) A reversible localizer other than $\operatorname{Fl}(\mathrm{Cat})$ contains $W_{H\mathbb Z}$, contains no truncation localizer $\mathbf W_n$, $n\geq0$ (a sphere $S^{n+1}$ is $\mathbf W_n$-aspherical), and never has enough intrinsic fibrations. It is moreover contained in $W_0$, that is, it satisfies Loc (5); this uses the description of the fundamental localizers not contained in $W_0$ and is not reproduced here.
(2) The fibration axiom Loc (7), in the form $D_{W_\omega}$, holds only for $\mathbf W_\infty$ and $\operatorname{Fl}(\mathrm{Cat})$ (commentary at ms.~65), and $\mathbf W_\infty$ is not reversible. So $\operatorname{Fl}(\mathrm{Cat})$ is the only localizer satisfying Loc (7) and reverse excision, and with Loc (5) there is none. The two groups of axioms of this chapter thus describe two different theories: ordinary homotopy types on the side of fibrations, and on the side of reversible cofibrations a family of localizers that begins at integral homology.
\end{dcomment}

% Chapter XVI section 4: working commentary by GPT-6 Astra.

\begin{dcomment}{exposition}{p28}{Three meanings of fibration}
Grothendieck is looking for maps along which an ordinary pullback already performs the correct homotopical operation. This explains the change of terminology at the beginning of the section: the class now called $\Fib_W$ is defined by the operation it makes possible, whereas the former fibrations, now $\Fibf_W$, are recognized by tests over contractible categories. The difference is the central issue of the next forty manuscript pages.

For clarity, write $W_S$ for the maps over $S$ whose underlying functor belongs to $W$, and write $W_S^{\mathrm u}$ for maps which remain in $W$ after every change of base $T\to S$. If $p:X\to S$, its intrinsic fibration property says that every $W$-equivalence $T'\to T$ over $S$ induces a $W$-equivalence
\[
 X\times_ST'\longrightarrow X\times_ST.
\]
Both $T$ and $T'$ vary. Testing only maps $T'\to S$ misses the required stability after a first change of base. By contrast, weak fibrancy makes the same demand only when $T$ and $T'$ are contractible in the categorical homotopy sense used in the preceding chapters. A category with contractible nerve need not admit a contraction through a finite chain of natural transformations; these two meanings should not be interchanged in reading the test.

The third notion, a Grothendieck fibration, concerns Cartesian lifts of arrows of $S$. A cleavage gives inverse-image functors $\alpha^*:X_t\to X_s$ for $\alpha:s\to t$. Membership in $\mathcal F_W$ means that these functors belong to $W$. Thus $\mathcal F_W$ combines strict categorical transport with homotopical local constancy. It supplies convenient representatives of families, but its definition does not say that pullback preserves every equivalence of bases.

Perfect maps enter from a different direction. Smoothness controls one comma-category comparison and properness the opposite comparison; perfection imposes both. The comparison with strict bifibrations was already discussed in the commentary on Chapter~VI, manuscript p.~186. Here perfection matters because it supplies replacements, and because the class of perfect maps is closed under composition. Those facts make Brown's argument available. No lifting-property definition of a Quillen fibration is being introduced.

The inclusion diagram is consequently asymmetric in a useful way. The geometric classes $\mathcal F_W$ and $\Parf_W$ sit inside the class detected on contractible bases; the intrinsic class $\Fib_W$ also sits there, but its inclusion in the other direction is a substantive assertion about $W$. Proving that reverse inclusion is the purpose of Loc~(7).
\end{dcomment}

\begin{dcomment}{exposition}{p31}{Why a reference localizer is retained}
The second localizer allows Grothendieck to separate the construction of a family from the equivalences which that family must preserve. A map may have transport belonging to a smaller class $W_0$, while the desired base-change theorem concerns the larger class $W$. This distinction is valuable precisely when changing the homotopy theory makes more maps into equivalences: additional equivalences of bases need not remain harmless for an existing family.

From $W_0\subset W$ one immediately obtains
\[
 \mathcal F_{W_0}\subset\mathcal F_W,\qquad
 \Fibf_{W_0}\subset\Fibf_W,\qquad
 \Parf_{W_0}\subset\Parf_W.
\]
Each defining condition on the left asks specified comparison functors to belong to the smaller class. There is no corresponding formal inclusion $\Fib_{W_0}\subset\Fib_W$. Here the quantified class of equivalences of bases has also enlarged. This is why the intrinsic fibrations appear at opposite ends of the diagram rather than in its vertical chain.

The pointwise pullback test at ms.~36 follows from this same distinction. For a class $\mathcal F$ stable under arbitrary base change, it suffices to check the outermost square with a map in $W$ at its bottom. Indeed, to test $p\in\mathcal F$ against a map $T'\to T$ over $S$, first pull $p$ back to $T$; the resulting map still belongs to $\mathcal F$, so the proposed single-square test applies. This proves the full intrinsic fibration property, with no assumption about the saturation of $W$ beyond the ordinary equivalence calculus.

The observation about retracts at ms.~29 deserves the same care as in the commentary on Chapter~XII, manuscript p.~209. If $W$ is closed under retracts, then $\Fib_W$ is closed under retracts: every pullback comparison for a retract of $p$ is a retract of the corresponding comparison for $p$. Without that hypothesis this argument does not establish the assertion. We retain the source's claim in the audit, with this dependency explicit, rather than silently replacing mild saturation by retract closure. It has no role in the composition and replacement arguments below.
\end{dcomment}

\begin{dcomment}{exposition}{p35}{Local constancy belongs to the coherent family}
Grothendieck's preferred representatives carry their coherence with them. A split Grothendieck fibration is described, in the variance used here, by a strict functor $S^{\mathrm{op}}\to\Cat$; forgetting the splitting leads to a pseudofunctor with canonically coherent transport. The display at ms.~33 prints $\underline{\mathrm{Hom}}(S,\Cat)$, whereas the displayed fibre functor at ms.~35 has the required opposite. We use the latter convention. The change is recorded as an edition slip and is not silently applied to his text.

The equivalence after localization between split and unsplit presentations does not mean that the forgetful functor is full before localization. A Cartesian functor may preserve chosen lifts only up to the unique vertical isomorphisms supplied by their universal property. Splitting replacements account for this discrepancy. The commentary on Chapter~VI, manuscript pp.~26 and~53--59, treats the comparison itself; its role here is to justify using either presentation when formulating the axiom.

There is also a distinction between saying that a coherent family is locally constant and retaining only its diagram in the ordinary homotopy category. The former is the object of $\HOT_W(S)$ together with the assertion that its transition maps are equivalences. Passing to $S^{\mathrm{op}}\to\Hot_W$ is a test for this property when $W$ is strongly saturated. It is not a reconstruction theorem for the family: higher compatibility data can disappear in the ordinary homotopy category. The commentary on Chapter~VI, manuscript pp.~7 and~73--74, gives the relevant examples; we do not repeat them here.

The modern comparison is especially transparent for ordinary homotopy types. \citet[Definitions~1.1--1.2, Theorem~1.3, and Corollary~1.28]{Cisinski2009LocallyConstant} characterize weak homotopy equivalences of indexing categories by their effect on locally constant homotopy-coherent diagrams. This is the later form of the geometric intuition behind Grothendieck's use of $\HOT_W^{\mathrm{lc}}$. The theorem concerns that coherent localization, with its coefficient category allowed to vary, and does not identify it with the category of ordinary functors into $\Hot$.

In the present argument the practical benefit is more elementary. The free split-fibration replacement $X\to\Phi_S(X)$ is natural in $X$, and for a weakly $W_0$-fibred $X\to S$ its comparison is universally a $W_0$-equivalence over $S$. Naturality will later let one replace not merely two objects but also the arrow between them. That extra feature is exactly what transfers condition B.
\end{dcomment}

\begin{dcomment}{later}{p36}{Intrinsic fibrations and sharp maps}
The definition chosen here has a later name which makes its independence of auxiliary lifting classes particularly clear. It is the definition of a sharp map in \citet[Section~2]{Rezk1998Sharp}: pullback along the map, after an arbitrary preliminary base change, preserves weak equivalences. In a right proper model category Rezk's Proposition~2.7 identifies these maps with those whose ordinary pullback squares are homotopy Cartesian. His Proposition~2.2 says that all model fibrations are sharp exactly when the model category is right proper.

These statements explain the intended geometry of $\Fib_W$, but they do not identify it with the fibrations of an arbitrary chosen model structure. Nor do they turn a weakly $W$-fibred category into an intrinsic fibration merely by changing terminology. The latter conclusion is Loc~(7), whose failure for connected-component equivalences is demonstrated below.

A useful consequence can already be read directly from the definition. Intrinsic fibrations are stable under composition: if $X\to Y\to S$ have this property and $T'\to T$ is a $W$-equivalence over $S$, first pull it back to $Y$, and then to $X$. Both comparison maps lie in $W$. They are also stable under every base change, because a test over the new base is already a test over the original one. These two elementary operations are the structural resources which the weaker class must acquire if it is to replace $\Fib_W$.
\end{dcomment}

\begin{dcomment}{exposition}{p39}{The three tests measure different operations}
Grothendieck separates three assertions which are easily conflated if one thinks only of familiar topological fibrations. Condition A concerns changing the base of an arbitrary $W$-equivalence. Condition B concerns changing the base of a $W$-equivalence between two suitable objects already lying over the same base. Condition C concerns one suitable map which is itself a $W$-equivalence. The implication from B to C uses the identity family $S\to S$ as the second object.

For a Grothendieck fibration in $\mathcal F_W$, the strict fibres are good representatives of the appropriate relative homotopy fibres. This accounts for the fibrewise reformulations of B and C. It does not remove their force: a total-space equivalence between locally constant families need not induce an equivalence on each strict fibre for a coarser homotopy theory. The covering obstruction at the end of the section will make this failure explicit in degree zero.

The added assertion F says only that every functor admits a factorization
\[
 X\xrightarrow{\ i\ }\bar X\xrightarrow{\ p\ }S,
 \qquad i\in W,\quad p\in\Fib_W.
\]
The word ``cofibrations'' in its printed definition is a slip for ``fibrations''. A separate dependency concerns the assertion that every object is intrinsically fibrant over the final category: this is exactly stability of $W$ under the functors $X\times-$. The editorial qualification to Rappel~2.3 must therefore be carried here; that product stability is established under Loc~(3 bis), not by the cited argument from the weak axioms alone. Factorizations by themselves should not silently be identified with an all-objects-fibrant Brown structure without this additional fact. There is no resulting gap in the implication from D: every $X\to e$ is perfect directly from its comma-category definition, so $\Parf_{W_0}\subset\Fib_W$ independently supplies all the required terminal fibrations, and therefore product stability. F still does not say that the already available class $\mathcal F_W$ consists of intrinsic fibrations.

The footnote correcting the recollection about $f_!$ is mathematically significant. Construction of the left homotopy Kan extensions is controlled here by relative Theorem~A, Loc~(3 bis); it should not be attributed to the much stronger B merely because B had accompanied earlier arguments. This is consistent with the later exposition of \citet[pp.~3--4]{Maltsiniotis2005Asphericite}. Smooth and proper base-change formulas, locally constant families, and existence of Kan extensions remain distinct pieces of the formalism.

The speculative independence from the cofibration axioms at ms.~37 is therefore a question about the scope of an abstract presentation, not a reason to distrust the familiar homotopy theory of spaces. Its present status must take account of the minimum-localizer theorem, discussed at ms.~65 below.
\end{dcomment}

\begin{dcomment}{examples}{p41}{An identity needed by the formal scheme}
Grothendieck now abstracts the three tests so that one proof can serve several classes of maps. The six classes to which he applies the scheme all contain the identities. That small hypothesis is needed for the assertion, made for an arbitrary $\mathcal F$, that B implies C.

Here is a complete test of the literal statement without it. Take $W=\mathbf W_\infty$ and let $\mathcal F$ consist only of the inclusion $j:e\to[1]$ selecting $0$. The full subcategory $\mathcal F_{[1]}$ has one object and only its identity endomorphism; for other bases it is empty. Thus every weak equivalence between its objects remains a weak equivalence after every base change, and $\mathrm B(\mathcal F)$ holds. Both $e$ and $[1]$ are contractible, so $j\in W$. Pulling $j$ back along the inclusion of $1$, however, gives $\varnothing\to e$, which is not a weak homotopy equivalence. Hence $\mathrm C(\mathcal F)$ fails.

The repair is simply to require $\id_S\in\mathcal F$ for every $S$. Apply B to the arrow $p:(X\to S)\to(S\to S)$ whenever $p\in\mathcal F\cap W$; the resulting universal equivalence is exactly C. All subsequent applications satisfy this hypothesis. This is a missing condition in the unrestricted formal sentence, not a defect in the six geometric cases.

A second notational repair keeps the scheme readable. In Lemma~4.4.1 the right factor is $p\in\mathcal F$, not the printed $f\in\mathcal F$. The literal latter requirement would demand that every map already belong to $\mathcal F$, which is neither used nor intended. The same distinction between the original arrow and its right factor recurs at ms.~50.
\end{dcomment}

\begin{dcomment}{exposition}{p43}{Why enough intrinsic fibrations gives B}
The first implication has a short proof which clarifies why factorization is relevant. Suppose every map admits a factorization as a $W$-equivalence followed by an intrinsic $W$-fibration. Let $u:X\to Y$ be a $W$-equivalence over $S$, where both structural maps $X\to S$ and $Y\to S$ are intrinsically fibrant. We must show that $u$ remains a $W$-equivalence after an arbitrary base change $t:T\to S$.

Factor that base change, rather than $u$, as
\[
 T\xrightarrow{\ a\ }Z\xrightarrow{\ b\ }S,
 \qquad a\in W,\quad b\in\Fib_W.
\]
Since $b$ is intrinsically fibrant, its pullback preserves $u$, so $u_Z:X\times_SZ\to Y\times_SZ$ lies in $W$. The structural maps of these two pullbacks over $Z$ are again intrinsically fibrant. Consequently the maps
\[
 X\times_ST\longrightarrow X\times_SZ,
 \qquad Y\times_ST\longrightarrow Y\times_SZ
\]
induced by $a$ both belong to $W$. In the commutative square formed by these maps and $u_Z$, two-out-of-three gives $u_T\in W$. The map $t$ was arbitrary, which proves $u\in W_S^{\mathrm u}$.

If $\mathcal F\subset\Fib_W$ and every map factors through $\mathcal F$ after a $W$-equivalence, the assumed intrinsic factorizations exist. The argument therefore proves $\mathrm A(\mathcal F)\Rightarrow\mathrm B(\Fib_W)$, and restriction gives $\mathrm B(\mathcal F)$. It also explains why the statement concerns the full subcategory of maps over $S$, rather than merely Cartesian arrows between fibred categories: no Cartesian condition on $u$ was used.

The remaining input in Corollary~4.4.2 is the availability of the six factorizations. Perfect replacement is inherited from the path-category construction in VII. Applying $\Phi_S$ to a weakly fibred representative gives the locally constant Grothendieck-fibration replacement. For the latter argument, the underbrace $q_0\in\mathcal F_{W_0}$ at ms.~42 should be read with $W$ in place of $W_0$, unless the entire argument has first been carried out relative to $W_0$. Both versions are available; their parameters should not be mixed within one factorization.

This completes the ``formal well-known argument'' suppressed in the source. It uses base change, two-out-of-three, and the specified factorizations. No claim about full faithfulness of a localization functor, and no strong saturation assumption, is needed.
\end{dcomment}

\begin{dcomment}{context}{p45}{The precise contribution of Brown's lemma}
Grothendieck's enthusiasm comes from a reversal of the expected direction of verification. To prove that a candidate fibration class preserves all weak equivalences under pullback, one may first verify that its trivial members remain trivial after pullback. If the class is closed under composition and base change, contains the maps to the final object, and supplies weak-equivalence--fibration factorizations, Brown's theory promotes this apparently smaller assertion to the full base-change theorem.

In the original paper, the factorization lemma is in Part~I, \S1, pp.~421--422, and the relevant pullback theorem is Part~I, \S4, Lemma~2, pp.~428--429 \citep{Brown1973}. Lemma~1 of that section first proves invariance of base change between fibrant objects over a fixed base; Lemma~2 then treats a weak equivalence of bases. The full dual argument has been supplied in the commentary to XVI, ms.~20--22. Here it yields $\mathrm C(\mathcal F)\Rightarrow\mathrm A(\mathcal F)$ under the hypotheses of Lemma~4.4.3.

Perfect maps satisfy the composition requirement before Loc~(7) is imposed. The locally constant Grothendieck fibrations need not. This single difference explains why each triple for $\Parf_W$ or $\Parf_{W_0}$ immediately collapses to the condition denoted D, while the corresponding triples for $\mathcal F_W$ and $\mathcal F_{W_0}$ remain separated for the moment.

Footnote~28 has these classes interchanged: its unconditional collapse concerns perfect maps. It is not evidence for an unproved unconditional use of Brown's lemma on $\mathcal F_W$. The body of the argument, especially Corollaries~4.4.4--4.4.6, keeps the missing composition condition visible.
\end{dcomment}

\begin{dcomment}{examples}{p47}{A finite failure of composition}
The composition condition in Corollary~4.4.5 can be tested on a very small family. The commentary on Chapter~I, manuscript pp.~118 and~122, already examined failures of composition. The following presentation singles out the precise class $\mathcal F_W$ used here and needs only connected components.

Put $W=\mathbf W_0$, the class of functors inducing bijections on $\pi_0$, and let $G$ be a nontrivial finite group. Define a category $B$ with two objects $b_0,b_1$ by
\[
 \Hom_B(b_0,b_0)=G,\quad \Hom_B(b_1,b_1)=\{1\},\quad
 \Hom_B(b_0,b_1)=G,\quad \Hom_B(b_1,b_0)=\varnothing.
\]
Write $a_h:b_0\to b_1$ for the arrow labelled by $h$ and prescribe $a_hg=a_{hg}$. There is a Grothendieck fibration $q:B\to[1]$ with fibres $BG$ and $e$. Its inverse-image functor along $0\to1$ is $e\to BG$. Since both categories are connected, $q\in\mathcal F_{\mathbf W_0}$.

Let $E$ be a transitive $G$-set with more than one element. The functor $r:B\to BG$ sending $g$ and $a_g$ to $g$ pulls the corresponding locally constant discrete family on $BG$ back to a discrete fibration $p:A\to B$. Every transport map of this family is a bijection of sets. Thus $p$ lies in $\mathcal F_{\mathbf W_0}$, and in fact its transports are isomorphisms of discrete categories.

The composite $qp$ is again a Grothendieck fibration. Its fibre over $0$ is the action groupoid $E\mathbin{/\mkern-6mu/}G$, which is connected by transitivity; its fibre over $1$ is the discrete category $E$. The inverse-image functor between these fibres is the inclusion of the objects of the action groupoid. On connected components it is the map
\[
 E\longrightarrow E/G=\{*\},
\]
which is not bijective. Consequently $qp\notin\mathcal F_{\mathbf W_0}$, although both $p$ and $q$ belong to that class. Taking $G=C_2$ with its regular action gives a base $B$ with two objects and five arrows, and $A$ with four objects and ten arrows; all categories in the example are finite.

The mechanism is the loss of monodromy information when a transport functor is required to preserve only connected components. A second locally constant family can detect precisely the information the first equivalence discarded. Corollary~4.4.5 says that C together with composition closure prevents this phenomenon. It does not say that C can be used without checking composition.
\end{dcomment}

\begin{dcomment}{exposition}{p50}{How to read the rectangle of conditions}
The large diagram records a small number of operations. Moving down a column first specializes from arbitrary base equivalences to equivalences between families, then to a family equivalent to its base, and finally replaces weakly fibred families by perfect ones. Moving to the right restricts the admissible families from $W$ to $W_0$. The first three downward arrows ultimately depend on the natural comparison with $\Phi_S$, whereas the perfect-map triple has already been handled by Brown's lemma.

The last arrow, $\mathrm D_{W_0}\Rightarrow\mathrm F$, has an especially direct meaning. Given $f:X\to S$, take its path factorization
\[
 X\xrightarrow{\ i\ }\bar X\xrightarrow{\ p\ }S,
 \qquad i\in W_0\subset W,\quad p\in\Parf_{W_0}.
\]
The condition $\mathrm D_{W_0}$, in its A-form, says that $p$ is an intrinsic $W$-fibration. This is the required factorization for F. The isolated display $\Parf_{W_0}\subset W$ in the source is therefore an edition slip for $\Parf_{W_0}\subset\Fib_W$; it is not an additional theorem asserting that all perfect maps are equivalences. The subsequent occurrences of $f$ as the right factor should likewise be $p$.

This interpretation also explains the role of the smallest reference localizer. It supplies the largest scope for the existence argument, because asking its perfect maps to preserve $W$ is the weakest of the displayed perfect-map conditions. The existence of replacements remains separate from the claim that every weakly $W$-fibred family is already a suitable replacement. The connected-component example will satisfy the first assertion and fail the second.

None of the solid arrows in this rectangle relies on the conditional arrows which return upwards. Those additional arguments enter only after the natural replacement lemma is established. Reading the diagram in that order avoids a circular use of the equivalence between $\mathcal F_{W_0}$ and $\Fibf_{W_0}$.
\end{dcomment}

\begin{dcomment}{exposition}{p54}{Why universal comparison maps transfer the tests}
Grothendieck's replacement lemma isolates the exact information needed from a comparison map. A global weak equivalence is insufficient because the tests concern all base changes. Suppose that every structural map $f:X\to S$ in $\mathcal F$ has a factorization
\[
 X\xrightarrow{\ i\ }\bar X\xrightarrow{\ p\ }S,
 \qquad i\in W_S^{\mathrm u},\quad p\in\mathcal G.
\]
The notation $W_S^{\mathrm u}$ concerns changes of the base $S$, not arbitrary changes of the codomain $\bar X$. That relative condition is exactly sufficient.

For A, let $T'\to T$ be a $W$-equivalence over $S$. The two comparison maps $i_{T'}$ and $i_T$ lie in $W$. If A holds for $\mathcal G$, the map $\bar X_{T'}\to\bar X_T$ also lies in $W$. Commutativity and two-out-of-three then put $X_{T'}\to X_T$ in $W$. This proves A for $\mathcal F$.

For C, suppose $f$ itself lies in $W$. Since $i\in W$, two-out-of-three gives $p\in W$. If C holds for $\mathcal G$, every map $p_T:\bar X_T\to T$ lies in $W$. Composing it with $i_T$ gives $f_T\in W$ for every $T\to S$. Thus $f$ is universally a $W$-equivalence. The change from $p$ to $g$ in the printed proof is only a change of name, but should not be mistaken for a second factorization.

B involves a further datum: an arrow between two families. Suppose a $W$-equivalence $u:X\to Y$ over $S$ fits into a commuting square
\[
\begin{tikzcd}
 X\arrow[r,"u"]\arrow[d,"i"']&Y\arrow[d,"j"]\\
 \bar X\arrow[r,"\bar u"']&\bar Y
\end{tikzcd}
\qquad
 i,j\in W_S^{\mathrm u},\quad \bar X,\bar Y\in\mathcal G_S.
\]
The square first implies $\bar u\in W$. Assuming B for $\mathcal G$, every pullback $\bar u_T$ is in $W$. The pulled-back square, with $i_T,j_T\in W$, then gives $u_T\in W$. Notice that the square is needed only for those $u$ which are in $W$; a replacement functor defined on all arrows is a convenient sufficient construction, not a logically necessary extra assumption.

The free split-fibration construction supplies such a functor. For weakly $W_0$-fibred objects, its unit lies in $W_{0,S}^{\mathrm u}\subset W_S^{\mathrm u}$, and the resulting structural map belongs to $\mathcal F_{W_0}$. Applying the three arguments with $\mathcal G=\mathcal F_{W_0}$ and $\mathcal F=\Fibf_{W_0}$ gives every announced equivalence. The reverse implications are restrictions along $\mathcal G\subset\mathcal F$.

This proof explains why it is legitimate to postpone the equivalences earlier in the section: they depend on a particular natural replacement, not on any assertion that all weakly fibred families are intrinsically fibrant. It also states the genuine limit of Lemma~4.4.9. Its objectwise hypotheses prove A and C; the text does not establish the B-transfer under those hypotheses alone. We leave that unrestricted assertion unclaimed.
\end{dcomment}

\begin{dcomment}{examples}{p57}{The two localizers have different jobs}
The infinite-path argument is designed to improve a family until both comma-category comparisons are controlled. Alternating the two one-sided constructions gives
\[
 X\longrightarrow\Phi(X)\longrightarrow\Psi\Phi(X)
 \longrightarrow\Phi\Psi\Phi(X)\longrightarrow\cdots.
\]
The odd and even subsequences are cofinal. Under sequential-colimit stability they show that the resulting object is both smooth and proper. If the input is weakly $W$-fibred, the successive comparison maps are universally $W$-equivalent over the base, and so is their colimit. This is the argument of VI, ms.~193--195 \citep{Grothendieck1990VI}; the commentary on Chapter~VI, manuscript p.~194, treats the alternation in detail.

Corollary~4.4.12 uses two independent conclusions from that construction. To transfer A from $\Parf_{W_0}$ to $\Fibf_W$, one needs a right factor in $\Parf_{W_0}$ and a comparison unit in $W_S^{\mathrm u}$. The underlying construction does not depend on a localizer. Its perfection uses Loc~(8 bis) for $W_0$, while the unit's universal $W$-equivalence uses Loc~(8 bis) for $W$ and weak $W$-fibrancy of the input. Lemma~4.4.9 then applies with exactly the required parameters.

The last sentence of the printed proof goes further, asserting the universal unit conclusion for $W_0$ as well. That conclusion is not available for an arbitrary weakly $W$-fibred input. The qualification is explicit in VI: for a general input the unit is a global equivalence, while the stronger relative conclusion requires weak fibrancy for the same localizer.

One can see the obstruction without a complicated path calculation. Let $G\ne1$ and take the Grothendieck fibration over $[1]$ classified by
\[
 F(0)=e,\qquad F(1)=BG,\qquad F(1)\longrightarrow F(0)
 \text{ the unique functor}.
\]
Its total category has an initial object, so its nerve is contractible. The map is weakly $\mathbf W_0$-fibred, since both fibres are connected, but it is not weakly $\mathbf W_\infty$-fibred. A perfect ordinary-homotopy replacement has contractible homotopy fibres over the contractible base: its total space is weakly equivalent to the original contractible total space. Consequently its fibre over $1$ is contractible, whereas the original strict fibre is $BG$. The unit on that fibre cannot be a weak homotopy equivalence. This rules out the extra $W_0=\mathbf W_\infty$ assertion while leaving the $W=\mathbf W_0$ assertion intact.

Thus the theorem is repaired by using the weaker unit statement it actually needs. No additional axiom and no change to the conclusion of Corollary~4.4.12 is necessary.
\end{dcomment}

\begin{dcomment}{exposition}{p61}{The summary theorem and its printed hypotheses}
The theorem gathers a substantial amount of reasoning, and its hypotheses should be read from the proved statements rather than from the abbreviated labels alone. With no extra hypotheses, there are the two downward chains, the restrictions to the reference localizer, and the implication from the weakest perfect-map condition to enough intrinsic fibrations. Sequential-colimit stability of $W_0$ collapses its column; stability for both localizers identifies the two columns. These are precisely Corollaries~4.4.11 and~4.4.12.

Two local errors in the summary can otherwise change its meaning. At the end of part~3, $\Fibf_W\subset\Fib_W$ is called Loc~(6 a); it is Loc~(7 a). In part~4(a), the printed assumption is Loc~(7 F) for $W_0$. The proof previously supplied in Corollary~4.4.6 assumes Loc~(7), which ensures that $\mathcal F_{W_0}$ is closed under composition. The diagram at ms.~63 also labels the upward arrow by $W_0(7)$. Footnote~30 does discuss the weaker F, but there the class is $\Fib_{W_0}$, not $\mathcal F_{W_0}$. Substituting the former class for the latter requires an argument which that footnote does not provide.

Accordingly, the direct proof in this section establishes part~4(a) under Loc~(7). This observation is a dependency correction, not a counterexample to the printed conclusion. A later argument using minimality settles the relevant comparison without that extra hypothesis, as explained next. It is better to distinguish these two proofs than to pretend that the summary's F follows from the preceding paragraphs.

Part~4(b) has its own independent mechanism. The previous proof in XII, Proposition~1, ms.~54--56, interprets a pullback comparison as the map on homology with locally constant coefficients. The equivalence of the locally constant coefficient categories supplies an isomorphism in $\Hot_W$; strong saturation then returns membership in $W$. Loc~(3 bis) supplies the homotopy Kan extensions used for the interpretation \citep{Grothendieck1990XIIa}. This explains both hypotheses, and why strong saturation cannot simply be replaced in that proof by closure under retracts.
\end{dcomment}

\begin{dcomment}{later}{p64}{Why ordinary homotopy types satisfy A}
The ordinary localizer provides a useful verification of the axioms, as well as the eventual answer to the minimum question. Start with a Grothendieck fibration $p:X\to S$ all of whose inverse-image functors are weak homotopy equivalences. For $s\in S$, the inclusion of the strict fibre $X_s$ into the comma category $s\backslash X$ has a right adjoint: it sends $(x,a:s\to p(x))$ to a Cartesian inverse image $a^*x$. The counit is the chosen Cartesian arrow. Hence this inclusion induces a homotopy equivalence on nerves. The transition functors between these comma categories induce, through these comparisons, the given inverse-image functors on strict fibres.

Quillen's theorem B, applied to the opposite categories, therefore identifies the strict fibres with the homotopy fibres of $Np$. A precise published formulation is \citet[paragraph~2.3.1 and Theorem~2.3.2, p.~136]{Cisinski2004Minimal}; the comma-category variance is why taking opposites is convenient here. If $p$ also belongs to $\mathbf W_\infty$, its homotopy fibres are weakly contractible. Thus every strict fibre $X_s$ is weakly contractible.

After any change of base $T\to S$, the resulting Grothendieck fibration has these same strict fibres. The dual theorem-A criterion now makes its projection a weak homotopy equivalence: the appropriate comma categories retract by adjunction onto the fibres, so they too are weakly contractible. We have proved C for $\mathbf W_\infty$, including its universal conclusion, without first assuming A. Empty bases cause no exception: a category mapping to the empty category is itself empty.

To recover A from this verification, the already proved Corollary~4.4.11 applies because $\mathbf W_\infty$ is stable under sequential filtered colimits. This stability can be seen without an additional axiom. The nerve commutes with filtered colimits: an $n$-simplex is a finite string of composable arrows, and the equations defining such a string are finite. The simplicial fibrant replacement $\operatorname{Ex}^{\infty}$ also commutes with filtered colimits, since each simplex of $\operatorname{Ex}^kK$ is specified by a map from the finite simplicial set $\operatorname{Sd}^k\Delta^n$, and its sequential union commutes with filtered colimits. A filtered colimit of Kan complexes is Kan because each horn has only finitely many simplices. Its components, its based homotopy classes of spheres, and the identifications given by homotopies are consequently represented at finite stages. A filtered diagram of weak equivalences therefore induces a bijection on components and isomorphisms on all homotopy groups in its colimit. Applying this to the nerves proves the required closure.

This gives the converse needed in the following status discussion: ordinary weak homotopy equivalences satisfy C, hence D, and then A by the infinite-path comparison. The use of filtered-colimit stability here proves a property of the ordinary localizer; it does not assert that an arbitrary weak fundamental localizer has that stability. Nor does the argument identify every categorical functor with a fibration. It identifies the particular locally constant Grothendieck fibrations with maps whose pullbacks perform the correct homotopy operation, and the natural replacement lemma extends that conclusion to the chapter's weak fibrations.
\end{dcomment}

\begin{dcomment}{status}{p65}{The minimum settles the implication questions}
Grothendieck is testing whether sequential limits are genuinely needed or merely belong to the replacement method he has available. The proposed candidate $W_\omega$ turns out to answer that question in the opposite direction from his prediction. \citet[Theorem~6.1.18, pp.~242--243]{Cisinski2006} proves that ordinary weak homotopy equivalences form the smallest \emph{weak} fundamental localizer. The definition used there is given in paragraph~6.1.12, p.~240, and uses the nonrelative theorem-A condition. Thus the theorem applies to the minimum of this chapter:
\[
 W_\omega=\mathbf W_\infty.
\]
The distinction from the earlier minimum theorem for the stronger relative axiom is essential. The assertion that $W_\omega$ almost certainly fails Loc~(8 bis) is therefore false: ordinary weak homotopy equivalences are stable under filtered colimits, as verified above.

Combining weak minimality with the fibre-detection theorem in \S6 settles the whole rectangle, without imposing Loc~(8 bis) on the two arbitrary localizers. For weak fundamental localizers $W_0\subset W$, all eight conditions are equivalent, and they hold precisely in the following two cases:
\[
 W=\mathrm{Fl}(\Cat)
 \qquad\text{or}\qquad
 W=W_0=\mathbf W_\infty.
\]
Here is the deduction, with the role of connectedness made explicit. Each of the eight conditions implies $\mathrm D_{W_0}$, which implies $\mathrm D_{W_\omega}$ by monotonicity of the perfect classes. If $W$ is all arrows, every test is automatic. Assume otherwise. Then $\mathrm D_{W_\omega}$ itself forces Loc~(5); one need not assume the potentially problematic converse from mere nondegeneracy to Loc~(5).

Indeed, an empty $W$-aspherical category would put $\varnothing\to e$ in $W$. Theorem~A would then put $\varnothing\to X$ in $W$ for every $X$, and two-out-of-three would give all arrows. If a nonempty $W$-aspherical $A$ were disconnected, choose $a:e\to A$ and factor it as $e\to E\xrightarrow{p}A$, with the first map in $W_\omega$ and $p$ $W_\omega$-perfect. The category $E$ is connected because $e\to E$ lies in the tautological subclass $W_\omega\subset\mathbf W_\infty$. The map $p$ belongs to $W$ and is universally in $W$ by $\mathrm D_{W_\omega}$. Its strict fibre over a different component of $A$ is empty, again giving the contradiction $\varnothing\to e\in W$. Thus every $W$-aspherical category is connected and nonempty.

Now factor any $f\in W$ as a $W_\omega$-equivalence followed by a $W_\omega$-perfect $p$. Every strict fibre of $p$ is $W$-aspherical by D, hence connected and nonempty. The inclusion $W_\omega\subset\mathbf W_\infty$ makes $p$ ordinarily perfect, so it compares these fibres with the corresponding comma categories by ordinary weak equivalences. The connected-component version of Theorem~A therefore makes $\pi_0(p)$ bijective. The left factor also induces a bijection on components, so $\pi_0(f)$ is bijective. This proves Loc~(5) directly. The commentary at ms.~85 gives the full fibre-detection argument and this nondegenerate reduction.

Theorem~6.1 now gives $W\subset\mathbf W_\infty$; weak minimality gives the reverse containment. Finally $\mathbf W_\infty\subset W_0\subset W$ forces $W_0=\mathbf W_\infty$. Conversely ordinary weak homotopy equivalences satisfy A by the preceding theorem-B argument, so all eight tests hold in the second case too.

Consequently both questions $\mathrm D\Rightarrow\mathrm A$ and $\mathrm D_{W_0}\Rightarrow\mathrm A$ have affirmative answers, without sequential-colimit assumptions. The suggested example at ms.~37 satisfying Loc~(7) while failing ordinary excision cannot occur: the ordinary and the all-arrow localizers both satisfy the excision axiom. This is a deduction from later weak minimality and the present chapter's perfect-replacement arguments, not the proof Grothendieck had in hand. F remains strictly weaker, as the connected-component example shows. The proper coarse localizer remembering only emptiness is not an exception to the result: it fails D, and the argument does not treat nontriviality alone as connectedness.
\end{dcomment}

\begin{dcomment}{exposition}{p67}{Why an open cover gives the needed integral comparison}
The brief digression uses a geometric fact already isolated in XII: an actual open cover has a homotopically faithful categorical integral. The commentary on Chapter~XII, manuscript p.~149, treats the cover criterion separately from the stronger converse-excision assertions. That separation matters here, because Grothendieck is trying to avoid assuming the very cofibration axiom he hopes to recover.

Let $X=U_1\cup U_2$ be a cover by full open subcategories, in the convention that each $U_i$ is closed under predecessors, and put $U_0=U_1\cap U_2$. Let $J$ have arrows $0\to1$ and $0\to2$, and form the Grothendieck construction $E$ of the diagram $U_0\to U_1$, $U_0\to U_2$. There is a canonical augmentation $q:E\to X$. We can verify directly that it is in every weak fundamental localizer.

For an object $x\in X$, consider the comma category $E/x$. If $x\in U_1\setminus U_2$, every arrow $u\to x$ starts in $U_1$, since $U_1$ is open. In particular an occurrence labelled by $2$ must have $u\in U_0$. Forgetting the label defines $T:E/x\to U_1/x$, with section $I$ selecting label $1$. Define $H:E/x\to E/x$ by changing label $2$ to label $0$ and keeping labels $0,1$. This is well-defined because every label-$2$ occurrence now lies in the intersection; it sends an arrow from label $0$ to label $2$ to the corresponding arrow within label $0$. The arrows of the span give natural transformations $H\to\id_{E/x}$ and $H\to IT$, while $TI=\id$. Thus $T$ is a categorical homotopy equivalence. Since $U_1/x$ has terminal object $x$, the category $E/x$ is $W$-aspherical. The case $x\in U_2\setminus U_1$ is symmetric.

If $x\in U_0$, all objects occurring in $E/x$ lie in $U_0$, and the construction is the product of the span index $J$ with $U_0/x$. The first factor has initial object $0$ and the second has terminal object $x$. Their categorical contractions give a contraction of the product by a finite chain of natural transformations; no general product-stability assertion about $W$ is needed. Hence it is aspherical. The theorem-A axiom now gives $q\in W$.

The homotopy-colimit interpretation identifies the class of $E$ with the homotopy pushout of the span; relative Theorem~A supplies the needed Kan-extension formalism. Strong saturation permits one to read an isomorphism in the localized category back as membership in $W$. Combined with the comparison criterion of Proposition~3.6, this is the intended route from enough intrinsic cofibrations to Loc~(6 a). The source explicitly labels its verification provisional. The open-cover comparison itself is secure; extending its argument to arbitrary pushout squares without the indicated hypotheses would reintroduce the invalid converse-excision principle already audited in XII.
\end{dcomment}

\begin{dcomment}{examples}{p69}{Enough fibrations can be computed from components}
Grothendieck's counterexample is effective because its replacement can be calculated completely. Put $W=\mathbf W_0$, and let $f:X\to Y$ be any functor. Write $I=\pi_0X$, $J=\pi_0Y$, and $\phi=\pi_0f$. Define
\[
 \bar X=\coprod_{\alpha\in I}Y_{\phi(\alpha)},
\]
where $Y_j$ is the full connected component of $Y$ indexed by $j$. Send an object $x$ of $X$ to the copy of $f(x)$ indexed by the component of $x$; arrows remain in that copy. This defines $i:X\to\bar X$, and the projections of the copies define $p:\bar X\to Y$ with $f=pi$.

Each summand of $\bar X$ is connected and nonempty, so the summands are exactly its connected components. Thus $\pi_0(i)$ is the identity on $I$, and $i\in\mathbf W_0$. To verify the stronger property of $p$, let $T\to Y$ be arbitrary. The pullback is a disjoint union of the components of $T$, with a copy indexed by each $\alpha$ whose target component agrees with theirs. Consequently there is a natural identification
\[
 \pi_0(T\times_Y\bar X)
 \cong \pi_0T\times_J I.
\]
If $T'\to T$ is a $\mathbf W_0$-equivalence over $Y$, it gives a bijection on the first factor, compatible with the maps to $J$. Its pullback therefore gives a bijection on the displayed fibre products. Hence $p\in\Fib_{\mathbf W_0}$, proving F for all functors. The formula includes empty categories: if $I$ is empty then both the replacement and every pullback are empty.

This factorization keeps precisely the component data of the source and fills each selected target component with a whole copy of that component. It does not try to preserve a nontrivial local system already present in $X\to Y$. That explains why enough intrinsic fibrations can coexist with failure of A, B, C, and even $\mathrm D_{W_\omega}$.

For the latter failure, take a transitive action of a group $G$ on a finite set $E$ with $|E|>1$. The action-groupoid projection $E\mathbin{/\mkern-6mu/}G\to BG$ has connected source and target, and so belongs to $\mathbf W_0$. It is locally trivial with discrete fibre $E$, hence perfect for every fundamental localizer. Pullback to the unique object of $BG$ gives $E\to e$, which is not a $\mathbf W_0$-equivalence. Thus $\mathrm C(\Parf_{W_\omega})$ fails. This is the covering obstruction discussed in the commentary on Chapter~VI, manuscript p.~43, expressed in a form suited to the present factorization.
\end{dcomment}

\begin{dcomment}{development}{p70}{The component-monodromy criterion}
Grothendieck's counterexample suggests a more precise question: which locally constant categorical families really are intrinsic fibrations for connected-component equivalences? A complete answer can be given inside the class of Grothendieck fibrations. It explains both the covering obstruction and the unusually simple replacement just constructed.

\textbf{Component-monodromy criterion.}\ours\quad
Let $p:X\to S$ be a Grothendieck fibration whose transport functors induce bijections on connected components. Let
\[
 L:S^{\mathrm{op}}\longrightarrow\Ens,\qquad L(s)=\pi_0(X_s),
\]
be its component functor. For each connected component $C$ of $S$, put $I_C=\operatorname{colim}_{C^{\mathrm{op}}}L$. Then the following conditions are equivalent: $p\in\Fib_{\mathbf W_0}$; every canonical map $L(s)\to I_C$, for $s\in C$, is a bijection; the monodromy action of the fundamental groupoid of $C$ on $L$ is trivial. Here trivial means that every loop acts identically, not that a chosen identification of the different fibres was part of the original data.

We first justify the component calculation on which the proof depends. For any Grothendieck fibration $Z\to T$, not necessarily locally constant, there is a canonical bijection
\[
 \pi_0Z\cong\operatorname{colim}_{T^{\mathrm{op}}}\pi_0(Z_t).
\]
A connected component on the left is generated by vertical arrows in the fibres and by arrows lying over arrows of $T$. Cartesian factorization writes every arrow as a vertical arrow followed by a Cartesian arrow. Its relation on components is therefore exactly the relation identifying a component of $Z_t$ with its inverse image in $Z_s$ along $s\to t$. These are precisely the relations in the set colimit. Conversely each such relation is represented by a Cartesian arrow. This proves both surjectivity and injectivity of the canonical map. The construction is natural in changes of base and in functors over the base.

Assume first that $p$ is intrinsically $\mathbf W_0$-fibrant. For any $s\in C$, the functor $e_s\to C$ is a $\mathbf W_0$-equivalence. Pull it back along $X|_C\to C$. The resulting functor $X_s\to X|_C$ induces a bijection on components. By the calculation above this is exactly $L(s)\to I_C$. Thus the second condition is necessary.

Conversely suppose all these canonical maps are bijections. They are natural in $s$, so they identify $L$ on each $C$ with the constant functor of value $I_C$. For any $t:T\to S$, the pullback fibration has component functor $L\circ t^{\mathrm{op}}$. Applying the component calculation component by component gives the natural formula
\[
 \pi_0(X\times_ST)
 \cong\coprod_{D\in\pi_0T} I_{\pi_0(t)(D)}.
\]
A $\mathbf W_0$-equivalence $u:T'\to T$ over $S$ bijects the indexing components and preserves their image in $\pi_0S$. On the corresponding summands its map is the identity of $I_C$. Hence it induces a bijection on the displayed component sets. This is the full intrinsic fibration test, proving the first condition.

It remains to identify the second condition with monodromy. All arrows of $L$ are bijections, so $L$ extends uniquely to the groupoid obtained from $S$ by formally inverting its arrows. Within a connected component, choose an object $s$. Every element of the colimit can be transported to $L(s)$ along a zigzag, and two elements of $L(s)$ have the same image precisely when they lie in the same orbit of the loop group at $s$. Thus $L(s)\to I_C$ is the orbit map. It is bijective precisely when every orbit has one element, equivalently when every loop acts trivially. This completes the proof, including components with empty fibres.

For $S=BG$, the criterion says exactly that a locally constant family with fibre component set $E$ is intrinsically $\mathbf W_0$-fibrant if and only if the $G$-action on $E$ is trivial. A transitive nontrivial action fails; a constant discrete family succeeds. The replacement at ms.~69 succeeds because its component functor is already the constant family indexed by the copies of each target component.

This is our explicit development of the source's example, not a priority claim. We checked the general smooth-functor treatment in \citet[Theorem~1.2]{Maltsiniotis2005Asphericite}, the sharp-map account in \citet[Section~2]{Rezk1998Sharp}, and the component and covering discussions in the commentary on the earlier chapters. Those sources supply the surrounding notions; we did not find this exact component-monodromy formulation there. The proof above is independent of any unverified identification with a particular model structure.
\end{dcomment}

% Chapter XVI, sections 5--7. Commentary by GPT-6 Astra.

\begin{dcomment}{exposition}{p71}{Two different ways of controlling size}
Grothendieck now asks whether a homotopy theory can be recovered from its small pieces. The two conditions he separates answer different parts of that question. Filtered-colimit stability says that a compatible family of already known equivalences remains an equivalence upon passage to the union. Accessibility says that equivalences have witnesses within a bounded range of sizes. Neither condition should be read as the assertion that a particular finite stage of an arbitrary convergent diagram is already an equivalence.

For example, let $A$ be a countable connected category. Its finite subcategories need not be connected, and a finite collection of arrows need not even generate a finite category: one endomorphism can generate infinitely many powers. An infinite cardinal bound behaves better. Closing at most $\kappa$ arrows under identities and finite composition still gives at most $\kappa$ arrows. This is why the canonical presentations below use an infinite cardinal and allow enlargement of a stage.

There is a terminological change between these pages and contemporary accessible-category theory. Here ``$\aleph_0$-accessible'' means generated, in the specified filtered-colimit sense, by categories or functors with at most countably many arrows. It does not mean generated by finitely presentable objects. In modern regular-cardinal terminology, a category with at most $\kappa$ arrows is $\kappa^+$-presentable. Indeed, a functor out of it is specified by at most $\kappa$ choices of objects and arrows, subject to at most $\kappa$ equations. In a $\kappa^+$-filtered colimit those choices and equations can be realized at a common stage. The same argument applies to commutative squares between functors, hence to the arrow category. This elementary verification is enough to translate the size estimates used here without silently shifting the cardinal index.

The path replacement serves a different purpose. It converts the question about an arbitrary arrow into a question about the fibres of one specially controlled arrow. The argument needs the replacement functor to commute with filtered colimits. It does not require its two output classes to be closed under those colimits. Grothendieck's parenthetical prediction on manuscript page~73 nevertheless contains an error: $W_\omega$ is in fact $\mathbf W_\infty$, and is filtered-colimit stable. We return to the precise minimality theorem after the proof of~6.1.
\end{dcomment}

\begin{dcomment}{exposition}{p74}{Taking a fibre after a filtered colimit}
Grothendieck wants to replace a closure condition on arrows by a closure condition on aspherical categories. Proposition~5.1 achieves exactly this, provided the path fibration is an intrinsic $W$-fibration. One small indexing choice makes the proof literal.

Write $p_i:E_i\to B_i$ for the perfect replacements and $p:E\to B$ for their filtered colimit. Choose $b\in B$. After replacing the indexing category by a cofinal filtered poset, choose a stage $i_0$ and $b_{i_0}\in B_{i_0}$ representing $b$. For $i\geq i_0$ let $b_i$ be its image. There is then a canonical isomorphism
\[
 E_b\cong\mathop{\rm colim}_{i\geq i_0}(E_i)_{b_i}.
\]
To see surjectivity on objects, represent an object of $E_b$ at a stage. The equality of its image with $b$ holds at a later stage, since equality in a filtered colimit of sets is eventually witnessed. The same argument applies to arrows, including the requirement that their image be the identity of $b_i$. Eventual equality also gives injectivity. Composition is inherited from a common later stage. This proves the displayed isomorphism directly; equivalently, filtered colimits commute with the finite-limit data defining categories and their strict fibres.

Suppose now that $W$-aspherical categories are closed under the relevant colimits. Since $f_i=p_i i_i$ and $i_i\in W_\omega\subset W$, each $f_i\in W$ implies $p_i\in W$. The hypothesis $D_{W_\omega}$ makes each $p_i$ universally a $W$-equivalence, so every fibre $(E_i)_{b_i}$ is $W$-aspherical. The displayed formula makes $E_b$ aspherical. The colimit replacement $p$ is again $W_\omega$-perfect because it is the canonical replacement of $\mathop{\rm colim}f_i$, not because arbitrary filtered colimits of perfect maps have been shown perfect. The aspherical-fibre criterion gives $p\in W$, and then $f\in W$ by two-out-of-three. Conversely, apply closure of $W$ to the maps $X_i\to e$. The proof for sequences is identical, using a tail of the sequence.

This is also the correction needed when a basepoint is chosen on manuscript page~82. A point of a colimit has a representative at some stage; it does not have an image in every earlier stage. Once a representative is fixed, the later images form the based diagram needed for iterated loops.
\end{dcomment}

\begin{dcomment}{examples}{p78}{Enlarging a countable subfunctor}
The order on pairs of subcategories is meant to describe simultaneous approximation of a functor, rather than separate approximation of its source and target. In the displayed formula on this page the right-hand side must be read as $X_i\subset X_j$ and $Y_i\subset Y_j$. The printed comparisons $X_i\subset Y_i$ and $X_j\subset Y_j$ do not even compare two stages. We retain that edition slip in the transcription.

Fix an infinite cardinal $\kappa$ and a functor $f:X\to Y$. Any collection of at most $\kappa$ pairs $(A_i,B_i)$ has an upper bound: take the subcategory of $X$ generated by all arrows of the $A_i$, and the subcategory of $Y$ generated by all arrows of the $B_i$ together with their necessary images under $f$. Each still has at most $\kappa$ arrows. A filtered union of at most $\kappa$ such pairs is itself a pair of the permitted size. Every arrow of $X$ or $Y$ occurs at some stage. These observations verify, without an appeal to abstract accessibility, the canonical adapted presentation used in Corollary~5.3.2.

Here is the admitted Lemma~5.6 in that presentation. Suppose $\pi_0(f)$ is bijective and begin with a countable subfunctor $f_0:A_0\to B_0$. We construct increasing countable subfunctors $f_m:A_m\to B_m$. At the next stage perform two operations. For every object $b\in B_m$, choose an object $a\in X$ whose image lies in the connected component of $b$, and choose a finite zigzag in $Y$ from $f(a)$ to $b$. Adjoin $a$, this zigzag, and their images to the relevant categories. Next, for every pair $a,a'\in A_m$ whose images are connected in $B_m$, choose a finite zigzag between $a$ and $a'$ in $X$. Such a zigzag exists by injectivity of $\pi_0(f)$. Adjoin it and its image. There are countably many choices to make, each with finite support. Closing under identities and composition preserves countability.

Let $A_\infty=\bigcup_m A_m$ and $B_\infty=\bigcup_m B_m$. Every object of $B_\infty$ occurs at some stage, so the first operation makes $\pi_0(A_\infty)\to\pi_0(B_\infty)$ surjective. If two objects of $A_\infty$ have connected images in $B_\infty$, the finite zigzag witnessing this lies at some stage containing both objects. The second operation at the next stage makes them connected in $A_\infty$. Thus the map is injective. We have found a countable enlargement belonging to $\mathbf W_0$.

Finally, $\pi_0$ preserves colimits because it is left adjoint to the discrete-category functor. A filtered colimit of bijections is a bijection, so $\mathbf W_0$ satisfies Loc~(8). The enlargement criterion now proves its asserted countable accessibility. Taking the target to be the terminal category gives the analogous result for connected categories. A small initial subcategory can therefore be enlarged to a connected one; it need not have been connected at the outset.
\end{dcomment}

\begin{dcomment}{exposition}{p79}{A direct completion of the accessibility argument}
Grothendieck's theorem concerns a concrete size bound, and one can prove it without completing the unfinished general statement of Lemma~5.5(b). We give that proof because it identifies exactly which witnesses must be added. The accessibility assertion of Theorem~5.4 is correct. Its parenthetical list of all axioms must, however, be qualified: the reverse cube condition Loc~(6b), if included in the bare label Loc~(6) as in Section~3, is not valid for ordinary weak homotopy equivalences. The valid direct excision assertion is Loc~(6a).

Put $R(A)=\operatorname{Ex}^{\infty}N(A)$. This is a Kan complex with the weak homotopy type of the nerve; the Kan construction and its proof are recalled in \citet[Proposition~2.1.39]{Cisinski2006}. Two elementary finiteness facts drive the argument. First, if $A$ has countably many arrows, then $R(A)$ has countably many simplices. In each degree, $\operatorname{Ex}$ is the set of maps from the finite simplicial set $\operatorname{Sd}\Delta^n$, and countably many finite choices remain countable. Iteration and the sequential union preserve that bound. Second, every finite simplicial configuration in $R(A)$ has support in a finitely generated subcategory of $A$. It appears at some finite $\operatorname{Ex}$ stage and involves finitely many original simplices; their finitely many arrows generate a countable subcategory. The same statement holds for a configuration together with a finite simplicial homotopy. These facts also show that $R$ commutes with filtered colimits.

Let $f:A\to B$ be in $\mathbf W_\infty$ and let $A_0\to B_0$ be a countable subfunctor. We build a sequence of countable enlargements. At stage $m$, first carry out the connected-component operations of the previous block. For each object $a\in A_m$, each $n\geq1$, and each element of
\[
 \pi_n(R(B_m),f(a)),
\]
choose a preimage in $\pi_n(R(A),a)$. Choose simplicial representatives and a homotopy exhibiting that the image of the chosen representative equals the prescribed class in $R(B)$. Adjoin the support of these choices. Next, for every element of $\pi_n(R(A_m),a)$ whose image is zero in $\pi_n(R(B_m),f(a))$, choose a null homotopy in $R(A)$ and adjoin its support. These choices exist because $R(f)$ is a weak equivalence. A countable Kan complex has countably many homotopy classes in every degree, so there are only countably many requirements over all $a$ and $n$. After adjoining images under $f$ and closing under composition, the next stage remains countable.

At the union, the connected-component argument still applies. To prove surjectivity of a positive homotopy group at an object $a$, represent a target class by a finite simplicial sphere in $R(B_\infty)$. Its support and the object $a$ occur at some stage. The next stage supplies the selected preimage and the required comparison homotopy. For injectivity, represent a class in $R(A_\infty)$ and a null homotopy of its image in $R(B_\infty)$. Both are finite configurations, so they occur together at some stage. The next stage supplies a null homotopy in the source. Consequently $A_\infty\to B_\infty$ is a weak equivalence. This proves the enlargement condition of Corollary~5.3.2 for $\kappa=\aleph_0$.

The remaining closure condition is equally finite. Given a filtered diagram of weak equivalences, a sphere, a null homotopy, or an equality of homotopy classes in its Kan replacement is represented at a stage. Bijectivity on components and homotopy groups therefore passes to the colimit. This proves Loc~(8) for $\mathbf W_\infty$ and completes the accessibility theorem. The construction is a direct proof of the known assertion of~5.4, rather than a claim of a new result.

The final reformulation of Lemma~5.5(b) also needs a qualification before it is used outside this example. The literal inclusion $\operatorname{Asph}_W\subset\Cat_{\rm connected}$ follows from Loc~(5); without it, the total localizer makes even the empty category aspherical. The unfinished extra hypothesis in the manuscript cannot be reconstructed uniquely, but this particular missing condition is visible.
\end{dcomment}

\begin{dcomment}{exposition}{p82}{Why the countable intersection is cofinal}
Grothendieck's loop argument eliminates the homotopy groups one at a time. The final passage from each degree separately to all degrees together is a small diagonal argument, and its closure hypothesis is essential.

Let $I$ be a directed poset in which every increasing countable sequence has a supremum. Suppose each $J_n\subset I$ is cofinal and closed under suprema of increasing countable sequences contained in $J_n$. Choose a sequence $n(0),n(1),\ldots$ in which every positive integer occurs infinitely often. Starting above an arbitrary $i\in I$, choose successively
\[
 i\leq i_0\leq i_1\leq\cdots,
 \qquad i_k\in J_{n(k)}.
\]
Cofinality makes every choice possible. Put $j=\sup_k i_k$. For a fixed $n$, the subsequence with $n(k)=n$ is cofinal in this increasing sequence. It has the same supremum $j$, and all its terms belong to $J_n$. Hence $j\in J_n$. This holds for every $n$, so $j\in\bigcap_nJ_n$. Since $j\geq i$, the intersection is cofinal.

Applied to the $I_n$ of the text, the argument finds a stage at which all iterated loop categories are connected. Connectedness of the original stage is already part of the reduction; together these conditions say that all homotopy groups vanish. The loop-detection comparison itself was treated in the commentary on Chapter~XII, manuscript p.~190. Here the new issue is making countably many conditions hold at one stage.

Two apparent defects should therefore be distinguished. The point $x$ must first be represented at a stage, as explained above. By contrast, the cardinal $\pi_0$ in the bound $\leq\pi_0$ is not a misprint: Theorem~5.4 has expressly used $\pi_0$ as another notation for $\aleph_0$. Replacing it is unnecessary and would erase the author's declared convention.
\end{dcomment}

\begin{dcomment}{exposition}{p85}{The fibration axiom recovers every homotopy group}
Grothendieck's comparison theorem explains why the seemingly modest axiom $D_{W_\omega}$ is so restrictive. Its role is not merely to furnish some factorization. It makes the fixed categorical path replacement valid for the chosen equivalences $W$, and that replacement retains all ordinary homotopy groups.

The proof has two logically separate steps. First, under $D_{W_\omega}$ a class $W$ is detected by its aspherical objects. If every $W$-aspherical category is $W'$-aspherical and $f\in W$, factor $f=pi$ with $i\in W_\omega$ and $p$ $W_\omega$-perfect. Then $p\in W$. The axiom makes all its fibres $W$-aspherical, hence $W'$-aspherical. Perfectness for $W_\omega$ implies perfectness for $W'$, and the aspherical-fibre criterion gives $p\in W'$. Since $i\in W'$, we have $f\in W'$. Only the source localizer $W$ needs the fibration axiom in this argument.

Second, let $X$ be $W$-aspherical and assume Loc~(5). It is connected and nonempty, so choose $x\in X$. The categorical based-path construction gives a map $E_xX\to X$ with $E_xX$ $W_\omega$-aspherical and with fibre the categorical loop object. Both total space and base are $W$-aspherical, so the map is in $W$. Since it is also $W_\omega$-perfect, its fibre is $W$-aspherical. Repeat at the constant loop. Loc~(5) says that every resulting loop category is connected. Thus all ordinary positive homotopy groups of $NX$ vanish. The first step, with $W'=\mathbf W_\infty$, now proves $W\subset\mathbf W_\infty$.

The connectedness hypothesis supplies the basepoint and turns asphericity into vanishing of components at every loop stage. The conclusion $W\subset\mathbf W_\infty$ cannot hold for the total localizer $W=\operatorname{Fl}(\Cat)$, which satisfies the universal-fibration condition but contains all functors. The shorter recollection on manuscript page~102 must therefore retain that exception or the explicit Loc~(5) hypothesis.

The total-localizer branch, cited rather than proved in the source, admits the following direct completion. Suppose that $W$ satisfies $D_{W_\omega}$ and that a $W$-aspherical category $A$ is disconnected. Choose an object $a$ in one component and an object $b$ in another. Since $A\to e$ is in $W$, two-out-of-three makes $a:e\to A$ a $W$-equivalence. Factor it canonically as
\[
 e\xrightarrow{i}E\xrightarrow{p}A,
 \qquad i\in W_\omega,\quad p\in\operatorname{Parf}_{W_\omega}.
\]
The tautological inclusion $W_\omega\subset\mathbf W_\infty$ makes $E$ weakly contractible, in particular nonempty and connected. This inclusion only uses that $\mathbf W_\infty$ is a fundamental localizer; it does not use the later equality with $W_\omega$. The image of $p$ is consequently contained in the component of $a$. Its strict fibre over $b$ is empty. On the other hand, $p\in W$ by two-out-of-three, so $D_{W_\omega}$ makes its empty fibre $W$-aspherical. Thus $\varnothing\to e$ belongs to $W$.

The same conclusion holds immediately if the original aspherical category is empty. Once $\varnothing$ is aspherical, Loc~(3) applies to $\varnothing\to B$ for every category $B$: each of its comma categories is empty and hence aspherical. Therefore $\varnothing\to B$ is in $W$ for all $B$. For any $f:B\to C$, two-out-of-three applied to $\varnothing\to B\to C$ gives $f\in W$. Hence $W$ is total.

Consequently, if $W$ is not total, every $W$-aspherical category is nonempty and connected. The loop argument above uses exactly this property: it applies to $X$ and successively to all its $W$-aspherical loop categories. It therefore proves $W\subset\mathbf W_\infty$ in the nontotal case, and in particular implies Loc~(5). This completes the whole dichotomy without a product-preservation or strong-saturation assumption. It also avoids the separate dependency issue with Proposition~2.7 discussed in the Section~2 audit; that general issue is not being used to withhold the stronger conclusion available under $D_{W_\omega}$.
\end{dcomment}

\begin{dcomment}{later}{p87}{The weak minimality theorem settles the finer-localizer hopes}
Grothendieck is looking for a homotopy theory finer than weak homotopy equivalence while retaining the elementary local axioms. The subsequent minimality theorem rules out that particular route. The exact result needed here is stronger than minimality among localizers satisfying the relative Theorem~A axiom.

\citet[Theorem~6.1.18, pp.~242--243]{Cisinski2006} proves that the weak homotopy equivalences of nerves form the least \emph{weak} fundamental localizer. His Definition~6.1.12 uses the ordinary aspherical-functor axiom, corresponding to this chapter's Loc~(3). Thus it applies to Definition~1.1 here, and gives
\[
 W_\omega=\mathbf W_\infty.
\]
The earlier strong-localizer theorem alone would not justify this substitution. The weak theorem does. The bracketed correction on manuscript page~87 is consequently valid in the precise axiom system of these pages.

Combining this equality with Theorem~6.1 gives a particularly transparent consequence. If $W$ satisfies Loc~(5) and $D_{W_\omega}$, then
\[
 \mathbf W_\infty=W_\omega\subset W\subset\mathbf W_\infty,
\]
so $W=\mathbf W_\infty$. There is no additional size or countable-colimit assumption in this deduction. With the completed dichotomy in the preceding block, $D_{W_\omega}$ alone allows exactly the two possibilities $W=\operatorname{Fl}(\Cat)$ and $W=\mathbf W_\infty$. In particular, the theory sought here cannot retain more information than ordinary homotopy types by choosing a smaller fundamental localizer.

The proposed $\widetilde{\mathbf W}_\infty$ requires a further reading decision. With Loc~(6) understood as the intended direct condition Loc~(6a), $\mathbf W_\infty$ satisfies the defining list, and the same sandwich makes the intersection equal to $\mathbf W_\infty$. Literally, however, Section~3 defines Loc~(6) to include the reverse condition~(6b). No fundamental localizer contained in $\mathbf W_\infty$ can satisfy that full list: weak minimality would make it $\mathbf W_\infty$ itself, which fails~(6b). Equivalently, Loc~(5) together with Loc~(7) forces $W=\mathbf W_\infty$, contradicting reverse excision. Thus the literal requested nondegenerate family is empty. One cannot call its purported minimum $\mathbf W_\infty$ without making the direct-excision qualification explicit. This is the same distinction required in the parenthesis of Theorem~5.4.

The theorem does not say that every weak fundamental localizer satisfies relative Loc~(3 bis). Minimality is an inclusion statement about the least class; relative locality is an additional property of an arbitrary larger class. Keeping these separate is necessary when reading the disputed recovery argument in~7.7.2.

There is a third reading, in which the bound $\widetilde{\mathbf W}_\infty\subset\mathbf W_\infty$ is dropped and one intersects all fundamental localizers satisfying the literal list, reverse condition (6b) included. By the dichotomy just recalled, Loc~(7) leaves only $\mathbf W_\infty$ and $\operatorname{Fl}(\mathrm{Cat})$, and $\mathbf W_\infty$ fails (6b). The intersection is then $\operatorname{Fl}(\mathrm{Cat})$. In each of the three readings the conjecture $\widetilde{\mathbf W}_\infty\subsetneq\mathbf W_\infty$ is false: the minimum is $\mathbf W_\infty$, or does not exist, or is the class of all functors.
\end{dcomment}

\begin{dcomment}{examples}{p90}{Why ordinary paths fail after truncation}
Grothendieck distinguishes a usable prescribed replacement from the bare existence asserted by Loc~(7F). A familiar example shows why that distinction survives even when a localized homotopy theory is completely well behaved.

Use the convention that $\mathbf W_n$ detects components and homotopy groups through degree $n$, and choose a connected based space $X=K(G,n+1)$ with $G\ne0$; for $n=0$ any nontrivial group is allowed, and for $n\geq1$ take an abelian group. The based path space $P_xX$ is contractible, so
\[
 p:P_xX\longrightarrow X
\]
is a $\mathbf W_n$-equivalence. Its ordinary fibre is $\Omega X$, whose homotopy group in degree $n$ is $G$ (whose components are $G$ if $n=0$). Hence $\Omega X\to e$ is not a $\mathbf W_n$-equivalence. Pulling $p$ back along the basepoint therefore destroys the equivalence.

The same calculation can be made with categorical models and the categorical path construction used in these pages: its nerves represent the indicated path space and loop space. It disproves universality of this particular $\mathbf W_n$-equivalence. It does not show that the localized theory has no homotopy fibre. Indeed, after $n$-truncation, both $P_xX$ and $X$ become terminal, so the homotopy fibre of the induced map in the localized theory is terminal. The nontrivial ordinary loop object was the wrong replacement for that calculation.

This explains the force of $D_{W_\omega}$. It insists that one path construction, which knows all homotopy groups, continue to calculate fibres after changing the equivalences. Truncation deliberately forgets a homotopy group of the base that becomes a visible homotopy group of the ordinary fibre. The obstruction is therefore structural. For $\mathbf W_0$, the explicit factorization at manuscript page~69 gives a useful replacement by copies of connected components; it avoids this obstruction and proves the weaker Loc~(7F). For every finite $n$, the proper localized Thomason models discussed at page~124 likewise prove Loc~(7F). Thus the source's caution about the usefulness of~F should be read as a concern about calculations, not as a theorem of nonexistence.
\end{dcomment}

\begin{dcomment}{development}{p95}{A sharp degree bound for the polygon examples}
Grothendieck's six-point example is decisive because subdivision doubles the number of edges without giving a categorical homotopy inverse. We can make its obstruction quantitative, and thereby see that the phenomenon persists through every iterated subdivision.

For $m\geq3$, let $C_m$ be the height-one poset with minimal elements $v_i$ and maximal elements $e_i$, indexed by $i\in\mathbf Z/m$, and with the only strict relations
\[
 v_i<e_i,\qquad v_{i+1}<e_i.
\]
It is the face poset of a polygonal circle. Its order complex is a circle with $2m$ edges. A choice of the cyclic ordering fixes the sign of the degree of a map of realizations.

\textbf{Degree spectrum for crowns.}\ours\quad
For $m,n\geq3$, the degrees of the maps $|NC_m|\to|NC_n|$ induced by monotone maps $C_m\to C_n$ are exactly the integers $d$ satisfying $|d|n\leq m$.

\emph{Proof.} A monotone map induces a simplicial map of the order complexes. Each of the $2m$ domain edges maps either to a single target edge or to a vertex. Give every target edge length one, and lift the image of the oriented domain circle to the universal covering line of the target circle, starting at a chosen lift of its initial vertex. Each traversed domain edge changes the lifted coordinate by at most one in absolute value. After all $2m$ edges the displacement is $2nd$, where $d$ is the degree. Therefore $2n|d|\leq2m$.

For the converse, choose a closed walk of $m$ steps in the set of polygon vertices of the target. Its first $n|d|$ steps go around the polygon $|d|$ times, in the sign prescribed by $d$; all remaining steps are stationary at the initial vertex. Send the $i$th minimal element of $C_m$ to the $i$th vertex of this walk. If a step traverses a target edge, send the corresponding maximal element to that edge. If the step is stationary, send the maximal element to the same target vertex. In either case the images of both adjacent minimal elements are less than or equal to the image of the maximal element. This defines a monotone map. Its realized walk has degree $d$, as required.\hfill$\square$

The nonempty-chain poset $\operatorname{sd}(C_m)$ is isomorphic to $C_{2m}$: its minimal elements are the singleton chains, and its maximal elements are the two-element chains, arranged in the subdivided cycle. Induction gives $\operatorname{sd}^r(C_m)\cong C_{2^r m}$. For every $r\geq1$ the proposition forces every monotone map
\[
 C_m\longrightarrow\operatorname{sd}^r(C_m)
\]
to have degree zero. Maps between circles of degree zero are null-homotopic, whereas the maximum map in the other direction has degree one. No categorical homotopy inverse can exist. Taking $m=3$ and $r=1$ proves the claim on manuscript page~95 in full.

There is also a direct categorical rigidity check. If an endomap $h:C_m\to C_m$ satisfies $h\geq\id$, it fixes the maximal elements. The image of $v_i$ must then lie below both $e_{i-1}$ and $e_i$ and above $v_i$, forcing $h(v_i)=v_i$. Thus $h=\id$. Dually $h\leq\id$ also forces equality. Any finite zigzag of natural transformations starting at the identity is therefore constant. A categorical homotopy equivalence between two crowns would have strictly inverse composites and hence be an isomorphism; crowns of different sizes cannot be categorically homotopy equivalent.

The general distinction between finite-space homotopy and simple homotopy is established in the literature discussed next. We did not find this explicit degree-spectrum calculation in the consulted finite-space papers. We include it as an elementary development of Grothendieck's example, without a priority claim.
\end{dcomment}

\begin{dcomment}{later}{p96}{The finite-poset version of simple homotopy}
Grothendieck's attempted localizer is ruled out by minimality, but his suggestion to work directly with finite ordered sets has a substantial later realization. These are two different claims, and the failure of the first does not diminish the second.

For a finite $T_0$ space, regarded as a finite poset, Barmak and Minian define elementary simplifications by deleting a weak point: a point whose strict lower set or strict upper set is contractible. Their Theorem~3.10 identifies the resulting simple-equivalence classes with simple homotopy types of finite simplicial complexes, through the order-complex and face-poset constructions. Their Theorem~4.13 gives the corresponding statement for maps \citep{BarmakMinian2008}. This is a precise version of the finite combinatorial program suggested here. It is not a fundamental localizer on all of $\Cat$ finer than $\mathbf W_\infty$.

The crown calculation explains why strict categorical homotopy equivalence is too fine for the desired comparison. A polygon and its subdivision have the same simple homotopy type, while their face posets have different strict categorical homotopy types. The elementary finite-space operations therefore must allow more than deletion of a point with a genuine order-theoretic deformation retraction. Conversely, passing all the way to weak homotopy equivalence forgets the simple structure. The successful theory sits between these two equivalence relations on finite models.

The invariant mentioned on manuscript page~88 also needs its classical input. For a homotopy equivalence $f:K\to L$ of connected finite CW complexes, Whitehead torsion is an element
\[
 \tau(f)\in\operatorname{Wh}(\pi_1L)
 =K_1(\mathbf Z[\pi_1L])/\langle\pm g\mid g\in\pi_1L\rangle.
\]
It vanishes exactly when the equivalence is simple, with the usual finite CW structures. It is attached to an equivalence and its simple structures, rather than being a canonical element attached to an unmarked homotopy type \citep[Section~7]{Milnor1966Whitehead}. Thus a proposed functorial invariant must specify its objects and morphisms before one asks whether it factors through a localization. The manuscript's suggestion of higher invariants does not specify enough data to identify it with a particular later torsion theory.

The page marked~94 is explicitly missing in the edition. The surrounding argument still supplies the finite counterexample and its purpose; we do not reconstruct an absent manuscript page from those conclusions.
\end{dcomment}

\begin{dcomment}{exposition}{p104}{A presentation supplies operations at every diagram shape}
The shift from $\Cat$ to a presheaf category is a change of computational model. Its purpose is to make the homotopy theory of diagrams available uniformly. A homotopy category alone contains too little information for this task: an object of $\operatorname{Ho}(\mathcal M)^S$ is a diagram whose relations commute only after localization, whereas an object of $\operatorname{Ho}(\mathcal M^S)$ retains a coherent realization of those relations. A derivator keeps the latter categories, their restrictions, and their homotopy Kan extensions. This distinction is why the existence of $f_*$ for every shape is stronger than the ability to write down a fibre for a single arrow.

Modern statements require some translation before they can be applied here. In \citet[Definition~1.5]{Groth2013}, Der~(1) concerns arbitrary small coproducts of shapes, Der~(2) says that evaluations jointly detect isomorphisms, Der~(3) gives both Kan adjoints, and Der~(4) calculates them pointwise. Grothendieck's original numbering separates the existence, calculation, and exactness assertions into Der~(4)--(6) and their starred partners; his Der~(3) is a contractibility condition. Moreover, the convention here uses presheaves $\mathcal M(S)=\mathcal M^{S^{\rm op}}$. Passing to opposite shapes translates this convention into the covariant diagrams used by many modern authors. Numbers and the words ``left'' and ``right'' cannot safely be transferred without that translation.

For a combinatorial model category the modern realization theorem is \citet[Proposition~1.30]{Groth2013}. The proof uses diagram model structures to derive restriction and ordinary Kan extension. It supplies the uniform operations sought in these pages. A complementary comparison is \citet[Theorem~3.33]{Lenz2018}: the homotopy prederivator of a cofibration category agrees with that of its quasicategory of frames. Its Corollary~4.1 then obtains the corresponding derivators from complete or cocomplete quasicategories. These are precise realization and comparison results. They do not say that every prederivator is representable by a model category, or that the original class $W$ can automatically be recovered from any equivalent presentation. The latter is a separate recognition argument, undertaken on manuscript pages~111--118.
\end{dcomment}

\begin{dcomment}{examples}{p109}{Finite coproducts of shapes and the length of a zigzag}
The assertion called quasi-tautological is correct in its original scope. Chapter~I, manuscript page~67, footnote~3, defines Der~(1) for a \emph{finite} sum of shapes. Let $(\mathcal C,U)$ and $(\mathcal E,V)$ be relative categories, with identities among their weak equivalences. There is a canonical equivalence
\[
 (\mathcal C\times\mathcal E)[(U\times V)^{-1}]
 \simeq \mathcal C[U^{-1}]\times\mathcal E[V^{-1}].
\]
Here is the mechanism. Every ordinary arrow $(a,b)$ factors as $(a,1)$ followed by $(1,b)$. A formal inverse of $(u,v)$ decomposes in the same fashion, using the inverses of $(u,1)$ and $(1,v)$. Horizontal and vertical arrows commute in every rectangle. The same commutation holds when one or both arrows are inverted, by multiplying the commutative square by the relevant inverses. Thus a word in the localized product can be rearranged into a word in the first coordinate followed by a word in the second. Conversely, two localized coordinate words define such a composite; the defining composition, identity, and inverse relations in either coordinate remain valid in the product localization. These operations are mutually inverse on hom-sets and preserve composition. This proves the equivalence, including its faithfulness.

Induction gives finite products. Since
\[
 \mathcal M^{(S\amalg T)^{\rm op}}
 \cong \mathcal M^{S^{\rm op}}\times\mathcal M^{T^{\rm op}},
\]
and objectwise weak equivalences correspond under this isomorphism, localizing proves finite Der~(1). The empty shape gives the terminal category as well. No localizer axiom is needed.

An explicit example explains why the same reasoning must not be extended to infinitely many factors without hypotheses. For $n\geq1$, let $J_n$ have objects $0,\ldots,n$, with precisely the nonidentity arrows between adjacent vertices directed from the even vertex to the odd vertex. Thus every nonidentity arrow has length one, and no two nonidentity arrows compose. Let $\mathcal M=\coprod_{n\geq1}J_n$ and declare every arrow a weak equivalence. In $\mathcal M^{\mathbb N}$ take objects $x$ and $y$ whose $n$th coordinates are the vertices $0$ and $n$ in $J_n$. Each coordinate has a finite zigzag joining these vertices, so their images are isomorphic in $\prod_n\operatorname{Ho}(\mathcal M)$.

There is nevertheless no arrow from $x$ to $y$ in $\operatorname{Ho}(\mathcal M^{\mathbb N})$. Such an arrow would be represented by a finite word of, say, $k$ ordinary arrows and formal inverses. In coordinate $n$, each letter changes the vertex by at most one. Hence the word can travel at most $k$ edges, whereas the endpoints in $J_n$ are $n$ edges apart. Taking $n>k$ gives a contradiction. The comparison to the infinite product is therefore not full. The obstruction is the absence of a uniform finite zigzag, not a defect in Grothendieck's finite axiom. Modern arbitrary-coproduct Der~(1) requires the additional homotopical structure supplied by the realization theorems. A related bounded-length obstruction already occurs in \citet[p.~307]{Thomason1980}, where the fundamental group of an infinite product of elementary categorical circles gives bounded integer families rather than all integer families.
\end{dcomment}

\begin{dcomment}{exposition}{p110}{What it means to realize a localized arrow}
The word ``realize'' in the proof of Der~(2) permits replacement of the source and target by weakly equivalent objects. It does not say that the localization functor is full on the original, fixed pair of objects. Even a formal inverse of a weak equivalence need not be represented by an ordinary arrow in the reverse direction between those same objects.

In a model category, the precise calculation uses a cofibrant replacement $QX\to X$ and a fibrant replacement $Y\to RY$. Every arrow $X\to Y$ in the homotopy category is represented by an ordinary map $QX\to RY$. Applied in a suitable diagram model category, this gives a map of replacement diagrams. If every component of the original localized arrow is an isomorphism, every component of the representing map is an isomorphism in the base homotopy category. Strong saturation of model weak equivalences makes those components weak equivalences. The representing map is therefore an objectwise weak equivalence, and becomes an isomorphism in the localized diagram category. The weak-equivalence replacements show that the original arrow is also invertible. This proves the pointwise conservativity required by Der~(2).

Grothendieck invokes the analogous Brown calculation for a cofibration or fibration category. In the finite-poset case, one constructs the necessary replacement diagram inductively, placing an object after its already treated predecessors and factoring the resulting latching map, or dually using matching maps. A functorial factorization can instead be applied directly to natural transformations; functoriality makes the chosen middle objects into a diagram. This is the specific advantage of the functorial constructions recalled on pages~99 and~102.

For all small shapes, a modern cofibration-category theorem carries more structure than a finite Brown theorem. \citet[Definition~1.5 and Warning~1.7]{Lenz2018} explicitly includes coproduct and sequential-colimit conditions in the ABC framework. The distinction matters: a proof that works by finite induction cannot simply be reindexed by an arbitrary small category. Conversely, a theorem proved for ABC cofibration categories should not be used to certify an unqualified claim about every finite Brown category. Lenz's Example~1.49 also exhibits a general relative category whose homotopy prederivator fails even weak strength. A class of denominators alone supplies a prederivator; the homotopical factorization and gluing conditions do substantive work.
\end{dcomment}

\begin{dcomment}{exposition}{p113}{Recovering the equivalences from constant diagrams}
The recognition calculation is a particularly useful part of this survey. It also records exactly where the later circular argument will enter. Let $\mathbb D=\mathrm{HOT}_W$, let $p_S:S\to e$, and write $c_S=p_S^*$. Assuming the left Kan extensions constructed under Loc~(3~bis), set
\[
 H(S,\xi)=p_{S!}c_S\xi.
\]
Adjunction identifies
\[
 \operatorname{Hom}_{\mathbb D(S)}(c_S\xi,c_S\eta)
 \cong
 \operatorname{Hom}_{\mathbb D(e)}(H(S,\xi),\eta).
\]
For $f:S\to T$, restriction on constant objects corresponds under this identification to precomposition with the canonical map
\[
 H(S,\xi)\longrightarrow H(T,\xi).
\]
Consequently restriction is fully faithful on constants precisely when these maps are isomorphisms for every $\xi$. For the nontrivial direction, bijectivity of the displayed hom-set maps for every $\eta$ implies the assertion by the ordinary Yoneda criterion for an isomorphism. This is the exact reason that a condition on all morphisms between constants can be tested by homology objects.

Write $q:\Cat\to\mathrm{Hot}_W$ for localization. Every object $\xi$ of this homotopy category is represented by a category $X$. The categorical Kan-extension calculation gives
\[
 H(S,qX)\cong q(X\times S),\qquad
 H(T,qX)\cong q(X\times T),
\]
with the comparison induced by $1_X\times f$. Taking $X=e$ shows that full faithfulness on constants forces $q(f)$ to be an isomorphism. Conversely, the product calculation identifies the comparison with $1_{qX}\times q(f)$, so invertibility of $q(f)$ implies all the required isomorphisms. The argument thus identifies the class of $\mathbb D$-equivalences with
\[
 \{f\mid q(f)\text{ is an isomorphism}\}.
\]
It identifies that class with $W$ itself only when $W$ is strongly saturated. This qualification is expressly present on page~111.

There are two different localization facts in this proof. The preceding block used finite products of \emph{categories being localized} to prove finite Der~(1). Here one needs the cartesian product operation on $\Cat$ to descend to its localization and to agree with the products used in the homology formula. That uses the localizer's stability properties and the preceding Kan-extension theory. Grothendieck's bracketed correction from Loc~(1,2,3) to Loc~(1,2,3~bis) is therefore retained in the audited reading; a formal finite-product localization lemma cannot remove the hypothesis from the whole recognition argument.

For $T=e$, a $W$-aspherical $S$ has $q(S)\cong e$, so the comparison is invertible for every $X$. Thus $c_S$ is fully faithful. This verifies the contractibility axiom in this setting without identifying a contractible diagram shape with an arbitrary diagram category: constant diagrams embed fully faithfully, while nonconstant diagrams can still contain substantial information.
\end{dcomment}

\begin{dcomment}{exposition}{p116}{Why the initial object computes this left adjoint}
There is no initial-versus-terminal mistake in the displayed calculation. Its direction records the presheaf convention. An object of $\mathcal M(S)$ is a functor $S^{\rm op}\to\mathcal M$. If $0$ is initial in $S$, it is terminal in $S^{\rm op}$, and hence
\[
 p_!\eta=\mathop{\rm colim}_{S^{\rm op}}\eta\cong\eta(0).
\]
For completeness, the cocone is given by the unique arrows $s\to0$ in $S^{\rm op}$. A cocone from $\eta$ to an object $z$ is uniquely determined by its component $\eta(0)\to z$, since compatibility forces every other component to be its composite with $\eta(s)\to\eta(0)$. This proves the universal property and the natural isomorphism.

Now take an object $\xi\in\mathcal M$, let $c\xi=p^*\xi$, and choose a cofibrant replacement of this diagram,
\[
 q:Q(c\xi)\longrightarrow c\xi,
\]
whose components belong to $\Sigma$. The derived counit of $\mathbf Lp_!\dashv p^*$ is represented by
\[
 p_!Q(c\xi)\longrightarrow p_!c\xi\longrightarrow\xi.
\]
Under the displayed evaluation isomorphisms this is simply $q_0:Q(c\xi)(0)\to\xi$, because the second arrow is the identity. It is a weak equivalence, so the derived counit is an isomorphism. A right adjoint is fully faithful exactly when its counit is an isomorphism. Therefore the constant-diagram functor is fully faithful, which is the required asphericity conclusion for this shape.

The argument proves a property of the derivator produced by the model structure. It has no need to identify its associated fundamental localizer with a previously given $W$. That is why Proposition~7.7.1 can be used to verify the original Der~(3) independently of the disputed recognition statement. The symbols $W$ and $\Sigma$ in the surrounding text should still be distinguished: $\Sigma$ denotes weak equivalences in the chosen model category, whereas $W_{\mathbb D}$ is a class of functors between small categories. Several printed subscripts switch them inadvertently; the correction ledger records those slips without rewriting the edition.
\end{dcomment}

\begin{dcomment}{status}{p118}{The circular inference and the separate recognition problem}
The two authorial footnotes are mathematically decisive. The proof that a suitable presheaf presentation yields a derivator concerns the prederivator $\mathrm{HOT}_W$. Proposition~7.7.1 then associates a strongly saturated fundamental localizer $W_{\mathbb D}$ satisfying the relative locality axiom to that derivator. To conclude the same relative locality assertion for the original $W$, one still needs
\[
 W=W_{\mathrm{HOT}_W}.
\]
The proof on pages~111--113 establishes this equality under strong saturation \emph{and} Loc~(3~bis). Applying it to prove Loc~(3~bis) therefore assumes the very conclusion under consideration. The derivator existence theorem does not remove that logical dependence. This is the circle explicitly marked by footnote~55 and acknowledged in footnote~56.

Strong saturation has a different route. If the presentation identifies the original weak equivalences with the inverse image of the model weak equivalences, the model category's saturation theorem supplies Loc~(1~ter). That statement does not require the invalid recognition inference. Thus the two conclusions in Corollary~7.7.2 should not receive a single undifferentiated verdict: the saturation argument has already been given on page~107, while the proposed proof of relative locality fails. The claim register records the latter argument as a sketch with a circularity gap. A circular proof does not by itself refute its conclusion, and no counterexample to that exact implication is asserted here.

Nor does weak minimality solve this separate problem. Theorem~6.1.18 of \citet{Cisinski2006} identifies the \emph{least} weak fundamental localizer. It does not say that every weak fundamental localizer satisfies the relative axiom. The distinction remains necessary throughout this commentary. The verified modern model-existence theorem near the end of the chapter assumes a strong basic localizer, so it cannot be used to supply the missing implication by changing terminology mid-proof.

There is also a broader recognition question: how much of a homotopy theory survives passage to all its diagram homotopy categories? A precise positive result is \citet[Theorem~37]{Arlin2020}, read in the February~2020 arXiv revision. A map of quasicategories whose induced prederivator morphism is an equivalence is itself an equivalence. The proof recovers equivalences on mapping spaces from diagram categories. This detects a given map; it does not assert that every abstract prederivator has a quasicategorical or model presentation. Together with Lenz's comparison theorem, it explains the strength of the derivator perspective while leaving the hypotheses of the particular equality $W=W_{\mathrm{HOT}_W}$ visible.
\end{dcomment}

\begin{dcomment}{exposition}{p119}{The detached base-change assertion}
The edition expressly marks six lines as out of context. Their symbols determine the square they refer to, but do not recover all of its missing hypotheses. The evident intended arrangement is
\[
\begin{tikzcd}
 X' \arrow[r,"q"] \arrow[d,"f'"'] & X \arrow[d,"f"]\\
 S' \arrow[r,"p"'] & S.
\end{tikzcd}
\]
The printed comparison is $p^*f_!\to f'_!q^*$, with the variance of these pages. The following explanation distinguishes two sufficient geometric situations: $f$ is $W$-smooth, or $p$ is $W$-proper. It invokes stability of smoothness under composition in the first case and the previously established ``forgetting'' statement on page~199 in the second. It would be unsafe to turn this fragment into a theorem for arbitrary commutative squares.

The structural meaning is clear. A base-change assertion says that restricting a homotopy Kan extension and extending the restricted diagram agree. It requires a square for which the two pointwise indexing constructions are homotopically equivalent. For comparison, \citet[Proposition~1.24]{Groth2013} gives an exact modern statement: a pullback square of categories is homotopy exact if its right vertical map is a Grothendieck fibration or its bottom horizontal map is a Grothendieck opfibration, in that article's covariant-diagram convention. The proof reduces the pointwise calculation to actual fibres using adjunctions. This illustrates the mechanism; it does not identify all $W$-smooth or $W$-proper functors with Grothendieck fibrations, nor does it reconstruct the absent surrounding sentence. The fragment and its editorial warning therefore remain in the text.
\end{dcomment}

\begin{dcomment}{later}{p123}{Both proposed minima give ordinary homotopy types}
The two minima asked for on pages~122--123 have an unequivocal answer. Let $W^{(a)}$ be the least weak fundamental localizer satisfying Loc~(6$'$a,8), and let $W^{(b)}$ be the least one satisfying Loc~(3~bis,6$'$a,8). The ordinary class $\mathbf W_\infty$ satisfies both lists: here the cube condition is explicitly the direct, weak form, so the reverse-excision problem of page~87 does not arise. It is therefore a member of each class being intersected. Conversely, weak minimality puts $\mathbf W_\infty$ in every weak fundamental localizer. Hence
\[
 \mathbf W_\infty\subset W^{(a)}\subset\mathbf W_\infty,
 \qquad
 \mathbf W_\infty\subset W^{(b)}\subset\mathbf W_\infty.
\]
Both equalities follow. By the accessibility theorem proved above, both minima satisfy Loc~(9). Thus part~(a), accessibility, is affirmative, whereas part~(b), strict inequality with ordinary weak homotopy equivalences, is negative. These answers use \citet[Theorem~6.1.18]{Cisinski2006}, not merely a theorem asserting minimality among already strong localizers.

The conclusion does not make other localizers irrelevant. Coarser localizers can distinguish different localizations of ordinary homotopy theory, and their derivators can retain different homotopy-limit and homotopy-colimit operations. What fails is this specific expectation of a new theory obtained by imposing the stated axioms below $\mathbf W_\infty$. The separate proposal to encode Whitehead torsion needs additional finite or relative structure, as in the finite-space construction discussed at page~96.
\end{dcomment}

\begin{dcomment}{later}{p124}{Generalized Thomason models and the role of properness}
The model-existence question has a substantial positive answer with precise hypotheses. \citet[Theorem~5.2.15, p.~225]{Cisinski2006} constructs, for every accessible \emph{strong} basic localizer $W$, a cofibrantly generated model structure on $\Cat$ with weak equivalences $W$ and Thomason cofibrations. It is left proper. It is right proper exactly when the localizer is proper in the sense used there. The reference is the published Ast\'erisque version; its pagination differs from the shorter author preprint. This theorem extends the classical structure recalled in the preceding Theorem~5.2.12. Its adjunction uses $\operatorname{cat}\operatorname{Sd}^2$ and $\operatorname{Ex}^2N$; the two subdivisions are part of the construction, rather than an identification of every categorical homotopy with a simplicial one.

This answers the generalized Thomason question in a large specified class. It does not prove that Grothendieck's weak axioms Loc~(6$'$a,8,9), without relative locality, imply that class of hypotheses. His phrase about extending Thomason's structure can naturally include its explicitly mentioned properness. With that reading, the ensuing conclusion about enough intrinsic cofibrations and fibrations is justified. If one only knows existence of an arbitrary Quillen structure, however, identifying its fibrations with the intrinsic class $\mathrm{Fib}_W$ requires an additional argument. The properness qualification explains exactly what makes the intended inference work.

The implication supplied by properness is elementary and worth making explicit. Suppose $(W,C,F)$ is a right proper model structure on a category with pullbacks, and let $p:E\to B$ belong to $F$. For any $B'\to B$, the pullback $p':E'\to B'$ is again a model fibration. If $a:A\to A'$ is a weak equivalence over $B'$, pull back $p'$ first to $A'$. Right properness applied to that fibration shows
\[
 E'\times_{B'}A\longrightarrow E'\times_{B'}A'
\]
is a weak equivalence. Thus every base change of $p$ preserves weak equivalences under pullback: $p$ is an intrinsic $W$-fibration. The model factorization of $f$ as a trivial cofibration followed by a fibration is therefore a factorization as a $W$-equivalence followed by an intrinsic $W$-fibration. This gives enough such fibrations, namely the content needed for Loc~(7F). Cisinski records the corresponding comparison in Example~6.1.14.

Dually, in a left proper model structure, a model cofibration remains one after cobase change, and pushout along it preserves weak equivalences. Hence it is an intrinsic $W$-cofibration. Factoring any map as a cofibration followed by a trivial fibration gives the required cofibration factorization. In the generalized Thomason theorem this conclusion is available because left properness is part of the theorem. The fibration conclusion requires the additional right properness hypothesis. One should therefore separate the affirmative realization theorem from the intrinsic-fibration question for an arbitrary localizer.

The finite-truncation examples specifically requested here do satisfy the stronger hypothesis. In \citet[Section~9.2.1, p.~364]{Cisinski2006}, $\mathbf W_n$ is generated by making the simplex category $\Delta/\partial\Delta^{n+2}$ aspherical. It is accessible and has aspherical generators; Corollary~4.3.28 proves such a localizer proper. Corollary~9.2.15, p.~369, identifies its simplicial weak equivalences with the maps inducing a bijection on components and isomorphisms on homotopy groups through degree~$n$, so this is the class meant in the manuscript, not merely a similarly named localization. Theorem~5.2.15 therefore gives a proper model on $\Cat$ for every finite $n\geq0$. Applying the preceding pullback argument proves that $\mathbf W_n$ has enough intrinsic fibrations. This answers the particular Loc~(7F) question on pages~91 and~124 affirmatively, even though $D_{W_\omega}$ fails for each of these strict truncations.

This distinction also resolves the apparent tension with the earlier truncated examples. A theory can possess valid derived right Kan extensions through a model presentation even though the fixed ordinary path replacement is invalid for its weak equivalences. Derived constructions use the fibrant replacements and lifting conditions of the localized model. They need not be the ordinary loop categories that retain every homotopy group. The explicit $\mathbf W_0$ factorization at manuscript page~69 supplies a concrete instance; the $K(G,n+1)$ example at page~90 explains why the unmodified ordinary path construction cannot provide a universal replacement for all truncations.

Outside the truncations the answer can be negative. Enough intrinsic fibrations force a localizer to be generated by its aspherical objects: under Loc~(7F) every $W$-equivalence lies in every fundamental localizer containing the functors $A\to e$ with $A$ $W$-aspherical (the argument is the middle step of the proof in the next commentary). For an accessible localizer this is in turn equivalent to properness, hence to right properness of its Thomason model structure; of that equivalence only the implication from aspherical generators to properness \citep[Corollaire~4.3.28]{Cisinski2006} is used in these pages, and the other steps were not checked for this commentary. The homology localizers fail the necessary condition, as the next commentary shows; so their Thomason model structures exist, and their fibrations are not intrinsic.
\end{dcomment}

\begin{dcomment}{development}{p124}{No homology localizer has enough intrinsic fibrations}
The question of this page for homological equivalences has a negative answer, for every generalized homology theory. For a spectrum $E$ let $W_E$ be the class of functors $f$ for which $E_*(Bf)$ is an isomorphism; for the Eilenberg--Mac~Lane spectrum of a ring $R$ this is the localizer $W_R$ of the $R$-homology equivalences. The statement was not found in the sources consulted; no priority is claimed.

\textbf{Theorem.}\ours
\textit{For every spectrum $E\neq0$ the fundamental localizer $W_E$ does not satisfy Loc (7 F): some functor has no factorization as a $W_E$-equivalence followed by an intrinsic $W_E$-fibration. In particular this holds for $W_R$, $R$ a nonzero ring.}

\textit{Proof.}
We apply the lemma and Proposition~2 of the commentary at ms.~27 (``Reverse excision excludes the fibration axioms''). The class $W_E$ is a strongly saturated fundamental localizer satisfying (3 bis) and containing $W_\infty$, and it is not $\operatorname{Fl}(\mathrm{Cat})$: $E_*(\emptyset)=0\neq E_*(e)$. The sphere $S^n$, $n\geq0$, is not $W_E$-aspherical, since $E_*(S^n)=E_*\oplus E_{*-n}$ and the projection to $E_*$ has the nonzero kernel $E_{*-n}$. So the lemma applies to $W_E$.

Let $F$ be the free group on $x,y$ and $\alpha(x)=x$, $\alpha(y)=yxyx^{-1}y^{-1}$. The space $BF$ is a wedge of two circles and $B\alpha$ induces the identity on $H_1$. Its suspension is a homology equivalence of the simply connected space $S^2\vee S^2$, hence a homotopy equivalence; by the suspension isomorphism $E_*(B\alpha)$ is an isomorphism. So $B\alpha\in W_E$. If $W_E$ satisfied (7 F), Proposition~2 would give $B\alpha\notin W_E$.\hfill$\square$

Three remarks. First, the localizers $W_E$ satisfy the cube rules (6a), (6b), (6c), by the long exact sequences of cofibre sequences; so the theorem is also an instance of the last assertion of Proposition~2. Second, $W_{H\mathbb Z}$ satisfies Loc (6a) and Loc (8) (homology commutes with filtered colimits), and it is accessible, Loc (9); the accessibility is standard for homology localizations and is not proved here. So these axioms do not imply Loc (7 F), which answers the question of ms.~105--106. Third, the truncation localizers $\mathbf W_n$ behave in the opposite way: they have enough intrinsic fibrations (preceding commentary) and fail reverse excision (commentary at ms.~17$'$).
\end{dcomment}
